\exercisegroup

\begin{enumerate}
\def\labelenumi{\arabic{enumi})}
\item
  设 X 是一个林德勒夫空间（TG，第 IX 章，附录 1，第 75 页，定义 1），且 X 是可解圈的。证明，对于每个点 \(x \in X\)，庞加莱群 \(\pi_1(X, x)\) 是可数的。
\item
  设 Y 是一致空间，X 是 Y 的子空间。设 \((\mathrm{U}_{i})_{i\in\mathrm{I}}\) 是 X 的一个开覆盖。
\end{enumerate}

\begin{enumerate}
\def\labelenumi{\alph{enumi})}
\tightlist
\item
  证明，存在 Y 中包含 X 的一个开邻域 V，以及 V 的一个开覆盖 \((V_{i})_{i\in I}\)，满足以下性质：
\end{enumerate}

\begin{enumerate}
\def\labelenumi{(\roman{enumi})}
\item
  对每个 \(i\)，有 \(\mathrm{U}_i = \mathrm{V}_i \cap \mathrm{X}\)；
\item
  对 I 的每个子集 J，若 \(\bigcap_{j\in J}V_j\neq \varnothing\)，则有 \(\bigcap_{j\in J}U_j\neq \varnothing .\)
\end{enumerate}

\begin{enumerate}
\def\labelenumi{\alph{enumi})}
\setcounter{enumi}{1}
\item
  假设 Y 是局部连通的，且各 \(U_{i}\) 连通。证明，存在满足这些条件的覆盖，使得各 \(V_{i}\) 是 Y 的连通开子集。
\item
  假设各 \(U_{i}\) 连通，并对每个 \(i \in I\) 取一点 \(x_{i} \in U_{i}\)。假设 \(\pi_{1}(U_{i}, x_{i})\) 在 \(\pi_{1}(X, x_{i})\) 中的像归结为单位元。证明，存在 Y 中包含 X 的一个邻域 V，使得对每个 \(x \in X\)，典范同态 \(\pi_{1}(X, x) \to \pi_{1}(V, x)\) 存在一个收缩映射。
\end{enumerate}

\begin{enumerate}
\def\labelenumi{\arabic{enumi})}
\setcounter{enumi}{2}
\tightlist
\item
  \begin{enumerate}
  \def\labelenumii{(\arabic{enumii})}
  \setcounter{enumii}{2}
  \tightlist
  \item
    设 \(X\) 是拓扑空间，且 \(x\) 是 \(X\) 的一个点。
  \end{enumerate}
\end{enumerate}

设 G 是一个群，且 \(\varphi\colon\pi_{1}(\mathrm{X},x)\to\mathrm{G}\) 是群同态。设 \((\mathrm{U}_{n})_{n\in\mathbf{N}}\) 是 X 的子空间的递减序列，其交集归结为点 \(x\)。对所有 \(n\)，令 \(\mathrm{G}_{n}\) 为 G 中这样的元素所成的集合：它们是在 \(\varphi\) 下某个包含于 \(\mathrm{U}_{n}\) 中的圈的道路类的像。

\begin{enumerate}
\def\labelenumi{\alph{enumi})}
\item
  若 G 可数，则序列 \((\mathrm{G}_{n})\) 稳定。（归结为对所有 n 都有 \((\mathrm{G}_{n}:\mathrm{G}_{n+1})\geqslant n+1\) 的情形，并取 G 的一个枚举 \((\gamma_{n})_{n}\)；递归构造元素 \(\lambda_{n}\in G_{n}\)，使得记 \(w_{n}=\lambda_{1}\ldots\lambda_{n}\) 时，有 \(\gamma_{i}\notin w_{n}G_{n+1}\)，其中 \(1\leqslant i\leqslant n\)。然后证明交集 \(\bigcap_{i}w_{i}G_{i+1}\) 非空。）
\item
  假设 G 是阿贝尔群，且其单位元是 G 中唯一的无限可除元素。则 \(\bigcap_{n} G_{n}\) 归结为单位元。（设 a 是交集中的一个元素；对每个 n，取 \(\gamma_{n} \in \pi_{1}(U_{n}, x)\)，使其在 \(\varphi\) 下的像等于 a。递归构造道路类 \(\lambda_{n} \in \pi_{1}(U_{n}, x)\)，使得 \(\lambda_{n} = \gamma_{n}\lambda_{n+1}^{d}\)，其中 d 是一个固定的整数且 \(\geqslant 2\)。证明 \(b = (d - 1)\varphi(\lambda_{1}) + a\) 是无限可除的。由此推出 a = 0。）
\item
  更一般地，若 G 是阿贝尔群，且 G 的每个无限可除元素都是挠元素，则 \(G_{n}\) 的交集由挠元素组成。
\end{enumerate}

(3) 本练习及下一练习的结果取自 J. W. CANNON 与 G. R. CONNER 的文章《关于一维空间的基本群》，载于《拓扑学及其应用》153 (2006)，第 2648–2672 页。

\begin{enumerate}
\def\labelenumi{\arabic{enumi})}
\setcounter{enumi}{3}
\tightlist
\item
  设 X 是连通且局部道路连通的拓扑空间。设 x 是 X 的一个点，并且 x 具有可数的邻域基。设 S 是一个集合，且 \(\varphi\colon\pi_{1}(\mathrm{X},x)\to\mathbf{Z}^{(\mathrm{S})}\) 是满射群同态。
\end{enumerate}

\begin{enumerate}
\def\labelenumi{\alph{enumi})}
\item
  证明 S 是可数的。
\item
  假设 X 是紧的；证明 S 是有限的。
\item
  假设 \(\pi_{1}(X,x)\) 是自由阿贝尔群；证明 X 是可解圈的。
\item
  假设 \(\pi_{1}(X,x)\) 是自由群；证明 X 是可解圈的。（利用自由群的下中心列各子群的交集归结为单位元这一事实。）
\end{enumerate}

\begin{enumerate}
\def\labelenumi{\arabic{enumi})}
\setcounter{enumi}{4}
\tightlist
\item
  设 P 是夏威夷耳环（III，第 336 页，练习 6；我们采用其中的记号）。对 \(n \in \mathbf{N}^*\)，令 \(\alpha_n \in \pi_1(\mathrm{P}, o)\) 为 \(\mathrm{C}_n\) 中一个圈的道路类，该道路类生成 \(\pi_1(\mathrm{C}_n, 0)\)。
\end{enumerate}

\begin{enumerate}
\def\labelenumi{\alph{enumi})}
\item
  证明映射 \(\mathrm{Hom}(\pi_{1}(\mathrm{P},0),\mathbf{Z})\to\mathbf{Z}^{\mathbf{N}}\)，\(\varphi\mapsto(\varphi(\alpha_{n}))_{n\in\mathbf{N}}\)，是单射群同态，且其像等于 \(\mathbf{Z}^{(\mathbf{N})}\)。（4）
\item
  群 \(\pi_{1}(P,0)\) 不是自由群；\(\pi_{1}(P,0)\) 对其导出群的商不是自由阿贝尔群。
\item
  设 A 是夏威夷群岛（III，第 338 页，练习 10）。证明从 \(\pi_1(A, o)\) 到 \(\mathbf{Z}\) 的每个同态都是平凡的。
\end{enumerate}

(4) 参见 B. DE SMIT，《夏威夷耳环的基本群不是自由群》，《国际代数与计算杂志》，第 2 卷，第 1 期 (1992)，第 33–37 页。

\exercisegroup

\begin{enumerate}
\def\labelenumi{\arabic{enumi})}
\item
  设 G 是拓扑群 O(2, R)。有 \(\pi_{0}(G) \simeq Z/2Z\) 和 \(\pi_{1}(G) \simeq Z\)；\(\pi_{0}(G)\) 的非平凡元作用于 \(\pi_{1}(G)\)，其映射为 \(t \mapsto -t\)。
\item
  记 H 为哈密顿四元数代数（TG，第 VIII 章，第 4 页），\((1, i, j, k)\) 为其典范基。记 N 为 H 的范数二次型；范数为 1 的四元数集是球面 \(S_{3}\)。令 \(H_{0}\) 为迹为零的四元数所成的实向量子空间，令 \(N_{0}\) 为 N 在 \(H_{0}\) 上的限制。
\end{enumerate}

\begin{enumerate}
\def\labelenumi{\alph{enumi})}
\item
  证明 TG，第 VIII 章，第 4 页，命题 4 中的同态使球面 \(\mathbf{S}_3\) 成为 \(\mathbf{SO}(\mathrm{N}_0)\) 的 2 次覆叠。
\item
  由此推出 SO(3, R) 的庞加莱群同构于 Z/2Z。具体描述 SO(3, R) 中 \(I_{3}\) 内的一个圈，其严格同伦类是该庞加莱群的唯一非平凡元。
\item
  设 \(\varphi\) 是从 \(\mathbf{S}_3 \times \mathbf{S}_3\) 到 SO(4, R) 的映射，由 \(\varphi(\mathbf{q}_1, \mathbf{q}_2)(\mathbf{x}) = \mathbf{q}_1 \mathbf{x} \mathbf{q}_2^{-1}\) 给出。证明 \(\varphi\) 使 \(\mathbf{S}_3 \times \mathbf{S}_3\) 成为 2 次覆叠。
\item
  推出 SO(4, R) 的庞加莱群同构于 Z/2Z。
\end{enumerate}

\begin{enumerate}
\def\labelenumi{\arabic{enumi})}
\setcounter{enumi}{2}
\tightlist
\item
  设 n 是自然数。令拓扑群 \(\mathrm{SO}(n+1,\mathbf{R})\) 作用于 \(S_{n}\)，即 \(R^{n+1}\) 的球面。
\end{enumerate}

\begin{enumerate}
\def\labelenumi{\alph{enumi})}
\item
  证明向量 \(e_{n+1}\) 的固定子可识别为拓扑群 \(\mathrm{SO}(n,\mathbf{R})\)。
\item
  设 \(f \colon \mathbf{I} \to \mathbf{S}_n\) 是连续映射。假设 \(f(t) \neq (0, \ldots, 0, 1)\) 对所有 \(t \in \mathbf{I} .\) 都成立。证明，存在连续映射 \(g \colon \mathbf{I} \to \mathrm{SO}(n + 1, \mathbf{R})\)，使得 \(f(t) = g(t) \cdot \mathbf{e}_{n+1}\) 对所有 \(t \in \mathbf{I} \)。
\item
  若 \(n \geqslant 3\)，证明 \(\mathrm{SO}(n, \mathbf{R})\) 到 \(\mathrm{SO}(n+1, \mathbf{R})\) 的包含映射在庞加莱群之间诱导同构。
\end{enumerate}

\begin{enumerate}
\def\labelenumi{\arabic{enumi})}
\setcounter{enumi}{3}
\tightlist
\item
  设 n 是自然数。令 \(B_{0}\) 为 \(\mathbf{SL}(n,\mathbf{R})\) 中由对角系数全为严格正数的上三角矩阵组成的子群。
\end{enumerate}

\begin{enumerate}
\def\labelenumi{\alph{enumi})}
\item
  证明从 \(\mathrm{SO}(n,\mathbf{R})\times\mathrm{B}_{0}\) 到 \(\mathbf{SL}(n,\mathbf{R})\) 的映射 \((g,u)\mapsto g\cdot u\) 是同胚（参见 INT，第 VII 章，第 91 页，命题 7）。
\item
  证明 SO(n, R) 到 SL(n, R) 的包含映射在庞加莱群之间诱导同构。
\end{enumerate}

\begin{enumerate}
\def\labelenumi{\arabic{enumi})}
\setcounter{enumi}{4}
\tightlist
\item
  设 G 是连通的可解圈拓扑群，设 \((\widetilde{\mathrm{G}}, p)\) 是 G 的万有覆叠；在 \(\widetilde{G}\) 上赋予拓扑群结构，使 p 成为群同态。记 Z 为 G 的中心，e 为其单位元，\(\widetilde{Z}\) 为 \(\widetilde{G}\) 的中心，\(\widetilde{e}\) 为其单位元。
\end{enumerate}

\begin{enumerate}
\def\labelenumi{\alph{enumi})}
\item
  设 \(f\) 是群 G 的连续自同构。证明，存在唯一的连续映射 \(\widetilde{f} \colon \widetilde{\mathrm{G}} \to \widetilde{\mathrm{G}}\)，使得 \(\widetilde{f}(\widetilde{e}) = \widetilde{e}\) 且 \(p \circ \widetilde{f} = f \circ p\)。证明 \(\widetilde{f}\) 是群自同构。证明映射 \(\varphi \colon f \mapsto \widetilde{f}\) 是从 \(\operatorname{Aut}(\mathbf{G})\) 到 \(\operatorname{Aut}(\widetilde{\mathbf{G}})\) 的单射群同态。
\item
  记 \(\operatorname{Out}(\mathbf{G})\) 为群 \(\operatorname{Aut}(\mathbf{G})\) 对内自同构子群的商；类似地定义 \(\operatorname{Out}(\widetilde{\mathbf{G}})\)。证明同态 \(\varphi\) 通过取商所诱导的同态 \(\psi\colon\operatorname{Out}(\mathbf{G})\to\operatorname{Out}(\widetilde{\mathbf{G}})\) 是单射群同态。
\end{enumerate}

\begin{enumerate}
\def\labelenumi{\arabic{enumi})}
\setcounter{enumi}{5}
\tightlist
\item
  设 G、H 是连通实李群，且 \(f: G \to H\) 是李群的满射态射。记 \(Z_{G}\) 为 G 的中心，\(Z_{H}\) 为 H 的中心。
\end{enumerate}

\begin{enumerate}
\def\labelenumi{\alph{enumi})}
\item
  证明以下断言等价：(i) 映射 f 使 G 成为 H 的覆叠；(ii) \(\operatorname{Ker}(f)\) 是离散的；(iii) f 通过取李代数诱导的态射 \(\mathrm{L}(f)\colon\mathrm{L}(\mathrm{G})\to\mathrm{L}(\mathrm{H})\) 是同构。
\item
  证明此时 \(f(\mathrm{Z}_{\mathrm{G}})=\mathrm{Z}_{\mathrm{H}}\)，且 \(f^{-1}(\mathrm{Z}_{\mathrm{H}})=\mathrm{Z}_{\mathrm{G}}\)。
\end{enumerate}

\begin{enumerate}
\def\labelenumi{\arabic{enumi})}
\setcounter{enumi}{6}
\tightlist
\item
  设 G 是实李群，记 \(G_{0}\) 为其单位分支，\(Z_{0}\) 为 \(G_{0}\) 的中心，\(\pi_{0}(G)=G/G_{0}\)。设 \((\widetilde{G}_{0},p_{0})\) 是 \(G_{0}\) 的万有覆叠，\(\widetilde{Z}_{0}\) 是 \(\widetilde{G}_{0}\) 的中心，\(\widetilde{e}\) 是其单位元。
\end{enumerate}

\begin{enumerate}
\def\labelenumi{\alph{enumi})}
\item
  设 \(\theta\colon\pi_{0}(\mathrm{G})\to\mathrm{Out}(\mathrm{G}_{0})\) 是与扩张 \(\mathcal{E}\colon\mathrm{G}_{0}\to\mathrm{G}\to\pi_{0}(\mathrm{G})\) 相关联的同态。证明文献所引处定义的上同调类 \(\omega(\pi_{0}(\mathrm{G}),\mathrm{Z}_{0},\theta)\)（它属于 \(\mathrm{H}^{3}(\pi_{0}(\mathrm{G}),\mathrm{Z}_{0})\)）为零。
\item
  假设存在李群 \(\widetilde{\mathrm{G}}\)、满射 \(p\colon\widetilde{\mathrm{G}}\to\mathrm{G}\)，以及同构 \(j\)，它把 \(\widetilde{\mathrm{G}}_{0}\) 同构到 \(\widetilde{\mathrm{G}}\) 的单位分支上，使得 \(p_{0}=p\circ j\)。证明 \(\omega(\pi_{0}(\mathrm{G}),\widetilde{\mathrm{Z}}_{0},\psi\circ\theta)=0\) 在 \(\mathrm{H}^{3}(\pi_{0}(\mathrm{G}),\widetilde{\mathrm{Z}}_{0})\) 中，其中 \(\psi\colon\mathrm{Out}(\mathrm{G}_{0})\to\mathrm{Out}(\widetilde{\mathrm{G}}_{0})\) 是练习 5 中定义的同态。
\item
  假设 \(\omega(\pi_{0}(\mathrm{G}),\widetilde{\mathrm{Z}}_{0},\psi\circ\theta)=0\)。设 \(\widetilde{E}\colon\widetilde{G}_{0}\to E\to\pi_{0}(G)\) 是以 \(\pi_{0}(G)\) 为商、以 \(\widetilde{G}_{0}\) 为核的扩张，且其相关联同态等于 \(\psi\circ\theta\)。证明 \(\mathrm{E}/\pi_{1}(\mathrm{G}_{0})\) 是以 \(\pi_{0}(G)\) 为商、以 \(G_{0}\) 为核的扩张，且其相关联同态等于 \(\theta\)。
\end{enumerate}

现在令 \(\gamma\) 为 \(\mathrm{H}^2 (\pi_0(\mathrm{G}),\mathrm{G}_0)\) 中满足 \(\gamma \cdot [\mathrm{E} / \pi_1(\mathrm{G}_0)] = [\mathrm{G}]\) 的唯一元素（所引文献，e)）。考虑同态 \(\delta \colon \mathrm{H}^2 (\pi_0(\mathrm{G}),\mathrm{Z}_0)\to\) \(\mathrm{H}^3 (\pi_0(\mathrm{G}),\pi_1(\mathrm{G}))\)，它由阿贝尔群正合序列 \(1\to \pi_1(\mathrm{G}_0)\to\) \(\widetilde{\mathbf{Z}}_0\to \mathbf{Z}_0\to 1\)（练习 6）关联。证明元素 \(\delta (\gamma)\)（属于 \(\mathrm{H}^3 (\pi_0(\mathrm{G}),\pi_1(\mathrm{G}))\)）不依赖于扩张 \(\widetilde{\mathcal{E}}\) 的选择；将其记为 \(\delta (\mathrm{G})\)。

d) 假设 \(\omega (\pi_0(\mathrm{G}),\widetilde{\mathrm{Z}}_0,\psi \circ \theta) = 0\)。存在李群 \(\widetilde{\mathrm{G}}\)、满射 \(p\colon \widetilde{\mathrm{G}}\to \mathrm{G}\) 以及同构 \(j\)，它把 \(\widetilde{\mathrm{G}}_0\) 同构到 \(\widetilde{G}\) 的单位分支上，使得 \(p_{0}=p\circ j\) 的必要且充分条件是 \(\delta(\mathrm{G})=0.\)（5）

(5) 例子参见 R. L. TAYLOR 的文章《非连通拓扑群的覆叠群》，《美国数学学会会刊》5 (1954)，第 753–768 页。

\begin{enumerate}
\def\labelenumi{\arabic{enumi})}
\setcounter{enumi}{7}
\tightlist
\item
  设 G 是李群，H 是 G 的闭子群；记 \(H_{0}\) 为 H 的单位分支，\(\pi_{0}(H)\) 为群 \(H/H_{0}\)。再记 p 为从 G 到 G/H 的典范满射。
\end{enumerate}

\begin{enumerate}
\def\labelenumi{\alph{enumi})}
\item
  映射 p 具有道路提升性质。
\item
  若 H 连通，则同态 \(p_{*}\colon\pi_{1}(\mathrm{G},e)\to\pi_{1}(\mathrm{G}/\mathrm{H},p(e))\) 是满射。
\item
  存在唯一映射 \(\varphi\colon\pi_{1}(\mathrm{G/H},p(e))\to\pi_{0}(\mathrm{H})\)，对于 G 中每条道路 \(\widetilde{c}\)，其起点为 \(e\) 且终点属于 H，它把 G/H 中圈 \(p\circ\widetilde{c}\) 的道路类对应到 \(\widetilde{c}(1)\) 模 \(\mathrm{H}_{0}\) 的类；它是群同态。
\item
  若 G 单连通，则 \(\varphi\) 是同构。特别地，H 连通当且仅当 G/H 单连通。
\end{enumerate}

\begin{enumerate}
\def\labelenumi{\arabic{enumi})}
\setcounter{enumi}{8}
\tightlist
\item
  设 G 是连通李群，\((\widetilde{\mathbf{G}}, q)\) 是 G 的万有覆叠，\(\widetilde{e}\) 是 \(\widetilde{G}\) 的单位元。设 H 是 G 的闭连通子群；记 \(i: H \to G\) 为典范注入，\(p: G \to G/H\) 为典范满射。
\end{enumerate}

\begin{enumerate}
\def\labelenumi{\alph{enumi})}
\item
  令 \(H_{1}\) 为 \(\overline{q}^{-1}(H)\) 的单位分支，令 \(q_{1}\colon\widetilde{\mathrm{G}}/\mathrm{H}_{1}\to\mathrm{G}/\mathrm{H}\) 为 q 通过取商诱导的映射。证明 \(q_{1}\) 使 \(\widetilde{G}/H_{1}\) 成为 G/H 的覆叠，且其在 \(p(e)\) 上的纤维同构于 \(\pi_{1}(\mathrm{G}/\mathrm{H},p(e))\)。
\item
  推出存在群的正合序列
\end{enumerate}

\[
\pi_ {1} (\mathrm{H} _ {1}, \widetilde {e}) \xrightarrow {q _ {*}} \pi_ {1} (\mathrm{H}, e) \xrightarrow {i _ {*}} \pi_ {1} (\mathrm{G}, e) \xrightarrow {p _ {*}} \pi_ {1} (\mathrm{G} / \mathrm{H}, p (e)).
\]

\begin{enumerate}
\def\labelenumi{\arabic{enumi})}
\setcounter{enumi}{9}
\tightlist
\item
  \begin{enumerate}
  \def\labelenumii{\alph{enumii})}
  \tightlist
  \item
    证明，对于每个整数 \(n \geqslant 2\)，群 \(\mathrm{SU}(n, \mathbf{C})\) 都是单连通的。（空间 \(\mathrm{SU}(2, \mathbf{C})\) 同胚于 \(S_{3}\)。令 \(\mathrm{SU}(n, \mathbf{C})\) 作用于 \(C^{n}\) 的单位球面，把点 \((0, \ldots, 0, 1)\) 的稳定子识别为群 \(\mathrm{SU}(n - 1, \mathbf{C})\)，并使用练习 9。）
  \end{enumerate}
\end{enumerate}

\begin{enumerate}
\def\labelenumi{\alph{enumi})}
\setcounter{enumi}{1}
\tightlist
\item
  证明，对于每个整数 \(n \geqslant 2\)，群 \(\mathbf{SL}(n, \mathbf{C})\) 都是单连通的。（令 \(B_{+}\) 为 \(\mathbf{SL}(n, \mathbf{C})\) 中由对角系数为实数 \textgreater{} 0 的上三角矩阵组成的子群。证明从 \(\mathrm{SU}(n, \mathbf{C}) \times \mathrm{B}_{+}\) 到 \(\mathbf{SL}(n, \mathbf{C})\) 的映射 \((g, u) \mapsto g \cdot u\) 是同胚。）
\end{enumerate}

\exercisegroup

\begin{enumerate}
\def\labelenumi{\arabic{enumi})}
\item
  设 X、Y 是拓扑空间，且 \(f: X \to Y\) 是连续映射。设 \((T, p)\) 是 Y-空间。X-空间上的典范下降数据 \(\tau\) 在 \(Z = X \times_{Y} T\) 上是有效的（IV，第 385 页，命题 1），典范映射 \(Z/R_{\tau} \to T\) 是单射且连续的。然而，它不一定是满射，也不一定严格。
\item
  设 Y 是由各空间 \(X_{i}\) 粘合得到的拓扑空间。令 X 为族 \((\mathrm{X}_{i})\) 的不交并所成的拓扑空间，令 \(f: X \to Y\) 为典范映射（TG，第 I 章，第 16 页）。对每个 i，置 \(Y_{i} = f(\mathrm{X}_{i})\)，并令 \(f_{i}: X_{i} \to Y_{i}\) 为 f 通过取子空间诱导的映射。映射 \(f_{i}\) 是双射且连续的。若映射 \(f_{i}\) 不是同胚，则相对于 f 的 X 上的典范下降数据不是有效的。例如参见 TG，第 I 章，第 94 页，练习 15。
\item
  令 B 为点 \((x,y)\) 在正方形 \([0,1]^{2}\) 中且满足 x=0、或 x=1、或 y=0、或 y=1、或存在 \(n\in N\) 使 nx=1 所成的集合。令 \(B_{1}\) 和 \(B_{2}\) 分别为 B 的子集，由 B 中的点 \((x,y)\) 组成，其中分别满足 \(y\leqslant\frac{1}{2}\) 和 \(y\geqslant\frac{1}{2}\)。
\end{enumerate}

构造一个不是覆叠的 B-空间 \((\mathrm{E}, p)\)，但使得 \(B_{1}\)-空间 \((\mathrm{E}_{\mathrm{B}_{1}}, p_{\mathrm{B}_{1}})\) 和 \(B_{2}\)-空间 \((\mathrm{E}_{\mathrm{B}_{2}}, p_{\mathrm{B}_{2}})\) 都是覆叠。

\begin{enumerate}
\def\labelenumi{\arabic{enumi})}
\setcounter{enumi}{3}
\tightlist
\item
  若拓扑空间的每个点都有由道路单连通邻域组成的邻域基，则称该空间局部道路单连通。
\end{enumerate}

\begin{enumerate}
\def\labelenumi{\alph{enumi})}
\item
  举出一个道路单连通但不是局部道路单连通的拓扑空间的例子。
\item
  设 X 是拓扑空间，G 是在 X 上适当作用的离散群。若 X 局部道路单连通，则 X/G 也是如此。
\end{enumerate}

\begin{enumerate}
\def\labelenumi{\arabic{enumi})}
\setcounter{enumi}{4}
\item
  设 G 在可解圈且完全正则的拓扑空间 X 上适当作用。置 Y = X/G，并记 \(f: X \to Y\) 为典范映射。证明典范群胚态射 \(\varpi(X)/G \to \text{Coeg}(f)\) 是同构。（使用 LIE，第 IX 章，§9，练习 17，改写 IV，第 403 页定理 3 的证明。）
\item
 设 X 是道路连通的拓扑空间，G 是连续作用于 X 上的群，且 x 是 X 中满足 \(g \cdot x = x\)（对所有 \(g \in G\)）的点。记 p 为从 X 到 X/G 的典范映射。设 H 是群胚，且 \(\varphi: \varpi(X) \to H\) 是群胚态射。假设对于每条道路 \(c\)，它在 \(\mathrm{X}\) 上，并且对于每个 \(g\in \mathrm{G}\)，都有 \(\varphi ([g\cdot c]) = \varphi ([c])\)，并且对于每个圈 \(c\in \Omega_x(\mathrm{X})\)，都有 \(\varphi ([c]) = e_{\varphi (x)}\)。
\end{enumerate}

设 \(c_{1}\) 和 \(c_{2}\) 是 X 中的道路，使得 \(p \circ c_{1}\) 与 \(p \circ c_{2}\) 具有相同的起点和相同的终点。证明 \(\varphi([c_{1}]) = \varphi([c_{2}])\)。

\begin{enumerate}
\def\labelenumi{\arabic{enumi})}
\setcounter{enumi}{6}
\tightlist
\item
  设 X 是一个可解圈拓扑空间，设 G 是一个在 X 上逆紧作用的离散群；记 \(m\colon G \times X \to X\) 为 G 的作用，\(\mathrm{pr}_2\colon G \times X \to X\) 为第二投影，\(f\colon X \to X / G\) 为典范满射。
\end{enumerate}

\begin{enumerate}
\def\labelenumi{\alph{enumi})}
\item
  证明群胚态射 \(\varpi(f)\) 是满射。
\item
  设 H 是一个群胚，设 \(\varphi\colon\varpi(X)\to H\) 是群胚态射，满足 \(\varphi\circ\varpi(m)=\varphi\circ\varpi(\mathrm{pr}_{2})\)。
\end{enumerate}

设 \(c_{1}\) 和 \(c_{2}\) 是 X 中的道路，使得道路 \(f \circ c_{1}\) 和 \(f \circ c_{2}\) 在 X/G 中严格同伦。证明 \(\varphi([c_{1}]) = \varphi([c_{2}])\)。

\begin{enumerate}
\def\labelenumi{\alph{enumi})}
\setcounter{enumi}{2}
\tightlist
\item
  证明存在唯一的群胚态射 \(\varphi'\)，从 \(\varpi(\mathrm{X/G})\) 到 H，使得 \(\varphi = \varphi' \circ \varpi(f)\)。由此得到 IV 第 403 页定理 3 的一个新证明。
\end{enumerate}

\exercisegroup

\begin{enumerate}
\def\labelenumi{\arabic{enumi})}
\tightlist
\item
  设 A 是 \(\mathbf{R}^2\) 的闭离散子集。
\end{enumerate}

\begin{enumerate}
\def\labelenumi{\alph{enumi})}
\item
  对于 R \textgreater{} 0，令 \(D_{R}\) 为这样的 \(z \in C\) 所成的集合，其中 \(|z| < R\)；证明 \(D_{R} - (A \cap D_{R})\) 的基本群同构于以 \(A \cap D_{R}\) 为基的自由群。
\item
  证明存在同胚 \(f\)，它把 \(\mathbf{R}^{2}\) 映入自身，并满足 \(f(\mathrm{A})\subset\mathbf{N}\times\{0\}\) 。
\item
  证明 \(\mathbf{R}^2 - \mathbf{A}\) 的基本群同构于 \(\mathrm{F(A)}\) 。
\end{enumerate}

\begin{enumerate}
\def\labelenumi{\arabic{enumi})}
\setcounter{enumi}{1}
\item
  设 A 是 \(R^{2}\) 中坐标为 \((1/n,0)\) 的点集，对 \(n\in N^{*}\) 而言。证明 \(R^{2}-A\) 的基本群具有连续统的基数。特别地，它不同构于群 F(A)。
\item
  设 \(((\mathrm{X}_{i}, x_{i}))_{i \in \mathrm{I}}\) 是带基点拓扑空间族；记 \((\mathrm{X}, x)\) 为该族的楔和（IV，第 421 页，例 4）。假设对每个 \(i \in I\)，点 \(x_{i}\) 在 \(X_{i}\) 中是闭点，并且具有邻域基 V，使得群 \(\pi_{1}(\mathrm{V}, x_{i})\) 归结为单位元。则典范同态 \(*_{i \in \mathrm{I}} \pi_{1}(\mathrm{X}_{i}, x_{i}) \to \pi_{1}(\mathrm{X}, x)\) 是同构。
\item
  设 P 是如 I，第 146 页练习 5 b) 中所定义的 \(R^{2}\) 子空间。令 a 为点 \((0,0,1)\)（它属于 \(R^{3}\)），令 C 为线段 \([a,(x,0)]\)（其中 \(x \in P\)）的并。
\end{enumerate}

\begin{enumerate}
\def\labelenumi{\alph{enumi})}
\item
  证明空间 C 局部道路连通且道路单连通。
\item
  设 D 是空间 C 与其对称像 -C 的并。证明空间 D 单连通，但不道路单连通。
\item
  证明庞加莱群 \(\pi_{1}(D,0)\) 的基数为连续统的基数。
\end{enumerate}

\begin{enumerate}
\def\labelenumi{\arabic{enumi})}
\setcounter{enumi}{4}
\tightlist
\item
  设 P 是 \(R^{2}\) 的子空间，如 I 第 146 页练习 5 b) 所定义；另记 \(P^{+}\)（分别为 \(P^{-}\)）为点 \((x,y)\in\mathrm{P}\) 的集合，其中 \(y\geqslant0\)（分别为 \(y\leqslant0\)）。设 a 是点 \((0,0,1)\)（位于 \(R^{3}\) 中），设 A 是线段 [a,0] 与各线段 \([a,(2r_{n},0)]\)（\(n\in N\)）的并。置 \(B^{+}=(P^{+}\times\{0\})\cup A\)，\(B^{-}=(P^{-}\times\{0\})\cup A\)，且 \(B=(P\times\{0\})\cup A\)。
\end{enumerate}

\begin{enumerate}
\def\labelenumi{\alph{enumi})}
\item
  证明空间 \(B^{+}\) 可解圈，且庞加莱群 \(\pi_{1}(B^{+}, a)\) 是以 N 为基的自由群。
\item
  证明 A 可缩到 a。
\item
  证明从 \(\pi_{1}(B^{+},a)*\pi_{1}(B^{-},a)\) 到 \(\pi_{1}(B,a)\) 的典范同态不是满射；它由 \(B^{+}\) 和 \(B^{-}\) 到 B 的典范包含诱导。
\end{enumerate}

\begin{enumerate}
\def\labelenumi{\arabic{enumi})}
\setcounter{enumi}{5}
\tightlist
\item
  设 X 是实直线 R，A 是 X 的子集。设 Y 是将 A 缩并为一点后由 X 得到的空间。
\end{enumerate}

\begin{enumerate}
\def\labelenumi{\alph{enumi})}
\item
  若 A 离散，则 Y 的庞加莱群是自由群。
\item
  若 A 有聚点，则 Y 的庞加莱群的基数为连续统的基数，且不是自由群。
\item
  若 A 是点集 \(1/n\)（\(n \in \mathbf{N}^*\)）的闭包，则空间 X/A 同胚于夏威夷耳环。
\end{enumerate}

\begin{enumerate}
\def\labelenumi{\arabic{enumi})}
\setcounter{enumi}{6}
\tightlist
\item
  设 G 是图，S 是其顶点集，A 是其定向边集。记 \(|G|\) 为其几何实现；将 S 识为 \(|G|\) 的子集，并记 \(p: I \times A \rightarrow |G|\) 为典范映射。
\end{enumerate}

\begin{enumerate}
\def\labelenumi{\alph{enumi})}
\item
  证明空间 \(|G|\) 紧（分别局部紧）当且仅当图 G 有限（分别局部有限）。
\item
  证明从 S 到 \(|\mathbf{G}|\) 的典范映射在 \(\pi_0(\mathbf{G})\) 与 \(\pi_0(|\mathbf{G}|)\) 之间诱导双射。
\item
  记 \(\varpi(|\mathrm{G}|)_\mathrm{S}\) 为 \(\varpi(|\mathrm{G}|)\) 的全子群胚，其顶点集为 S。证明存在唯一同构 \(\varpi (\mathrm{G})\to \varpi (|\mathrm{G}|)_{\mathrm{S}}\)，它在顶点集上为恒等，并将定向边 \(a\in \mathrm{A}\) 的道路类映为道路 \(t\mapsto p(t,a)\) 的类
\item
  设图 G 连通且有限；令 \(m = 1 + \text{Card}(S) - \text{Card}(A)\)。证明 \(|G|\) 的庞加莱群是有 m 个生成元的自由群。
\end{enumerate}

\begin{enumerate}
\def\labelenumi{\arabic{enumi})}
\setcounter{enumi}{7}
\tightlist
\item
  设 X 是空间 \(R \times \{0,1\}\) 的商空间，商所取的是使 \((x,0)\) 与 \((x,1)\) 对所有 \(x \in R^{*}\) 均等价的最粗等价关系（“原点加倍的实直线”）。
\end{enumerate}

\begin{enumerate}
\def\labelenumi{\alph{enumi})}
\item
  证明 \(\pi_1(\mathrm{X},x)\) 同构于 \(\mathbf{Z}\)，对每个点 \(x\in \mathrm{X}\) 均如此。
\item
  设 \(\mathrm{X}^{(2)}\) 是空间 \(X^{2}\) 在对称群 \(S_{2}\) 的作用下所得的商空间，其中该作用通过置换坐标于空间 \(X^{2}\) 上进行。证明 \(\pi_{1}(\mathrm{X}^{(2)},y)=0\)，对所有 \(y\in\mathrm{X}^{(2)}\) 均成立。
\end{enumerate}

\begin{enumerate}
\def\labelenumi{\arabic{enumi})}
\setcounter{enumi}{8}
\tightlist
\item
  设 X 是拓扑空间，设 R 是 X 上的等价关系，设 \(p: X \to X/R\) 是典范满射。设 \(x \in X\)。
\end{enumerate}

\begin{enumerate}
\def\labelenumi{\alph{enumi})}
\item
  设 X 连通且局部道路连通，等价类是连通集，且空间 X/R 可解圈。证明同态 \(p_{*}:\pi_{1}(X,x)\to\pi_{1}(X/R,p(x))\) 是满射。
\item
  设 X = I，且 \(\{0,1\}\) 是唯一不退化为单点的等价类。证明同态 \(p_{*}\) 不是满射。
\end{enumerate}

\begin{enumerate}
\def\labelenumi{\arabic{enumi})}
\setcounter{enumi}{9}
\tightlist
\item
  设 W 是华沙圆，定义见 IV 第 455 页练习 1，现重新采用其记号。
\end{enumerate}

\begin{enumerate}
\def\labelenumi{\alph{enumi})}
\item
  设 R 是 W 上的等价关系，其唯一的非单点等价类为 B 和 S；设 \(p: W \rightarrow W/R\) 为典范满射。证明商空间 W/R 同伦等价于圆但不同胚于圆 \(S_{1}\)。证明同态 \(p_{*}\) 不是满射。
\item
 设 R' 是 W 上的等价关系，使 B ∪ S 成为唯一不退化为单点的等价类；设 \(p'\) : W → W/R' 为典范满射。证明空间 W/R' 同胚于 \(S_{1}\)，同态 \(p_{*}\) 不是满射，但 R' 的等价类是连通的。
\end{enumerate}

\begin{enumerate}
\def\labelenumi{\arabic{enumi})}
\setcounter{enumi}{10}
\tightlist
\item
  设 X 是道路连通拓扑空间，设 \(f: X \to X\) 是同胚。设 T 是 \(X \times I\) 的商空间，商所取的是使点 \((x,0)\) 与 \((f(x),1)\) 对每个 \(x \in X\) 均等价的最细等价关系。记 \(p: X \times I \to T\) 为典范映射；记 \(j\colon X\to T\) 为由 \(x\mapsto p(x,1/2)\) 给出的映射。设 a 是 X 中的一点，设 \(c\in\Lambda_{f(a),a}(X)\) 是起点为 \(f(a)\)、终点为 \(a\) 的道路，并记 \(\gamma=[c]\)。
\end{enumerate}

\begin{enumerate}
\def\labelenumi{\alph{enumi})}
\item
  证明从 I 到 T 的映射由 \(t \mapsto p(a, \frac{1}{2} - t)\)（当 \(t \in [0, \frac{1}{2}]\)）以及 \(t \mapsto p(c(2t - 1), \frac{3}{2} - t)\)（当 \(t \in ]\frac{1}{2}, 1]\)）给出，是 T 中基于 \(j(a)\) 的圈。记其类为 \(\delta\)。
\item
 设 S 是从 \(\pi_{1}(\mathrm{X},a)*\mathbf{Z}\) 到 \(\pi_{1}(\mathrm{T},j(a))\) 的唯一群同态，它将类 \(v\in\pi_{1}(\mathrm{X},a)\) 映为 \(j_{*}(v)\)，并将 Z 中的元素 t=1 映为类 \(\delta\)。证明同态 S 是满射，且其核是包含元素 \(vt(\gamma^{-1}f_{*}(v)\gamma)^{-1}t^{-1}\)（\(v\in\pi_{1}(\mathrm{X},a)\)）的最小正规子群。
\item
  对 \(v \in \pi_{1}(X, a)\)，置 \(\varphi(v) = \gamma^{-1} f_{*}(v) \gamma\)；证明 \(\varphi\) 是群 \(\pi_{1}(X, a)\) 的自同构。设 \(\varphi: \mathbf{Z} \to \operatorname{Aut}(\pi_{1}(X, a))\) 是将 1 映为 \(\varphi\) 的唯一群同态。构造从群 \(\pi_{1}(T, j(a))\) 到 Z 与 \(\pi_{1}(X, a)\) 的外半直积（对应作用为 \(\varphi\)）的同构。
\end{enumerate}

\begin{enumerate}
\def\labelenumi{\arabic{enumi})}
\setcounter{enumi}{11}
\tightlist
\item
  \begin{enumerate}
  \def\labelenumii{\alph{enumii})}
  \tightlist
  \item
  设 D 是 \(R^{3}\) 中的一条直线，令 \(U = R^{3} - D\)。证明 U 道路连通，并且对每个点 \(a \in U\)，群 \(\pi_{1}(U, a)\) 同构于 Z。
  \end{enumerate}
\end{enumerate}

\begin{enumerate}
\def\labelenumi{\alph{enumi})}
\setcounter{enumi}{1}
\tightlist
\item
  设 P 是 \(R^{3}\) 中的一个平面，C 是包含于 P 的圆，令 \(V = R^{3} - C\)。证明 V 道路连通，且对每个点 \(a \in V\)，群 \(\pi_{1}(V, a)\) 同构于 Z。（将范坎彭定理应用于由平面 P 界定的两个半空间构成的 \(R^{3}\) 的闭覆盖；也可以注意到空间 U 和 V 同胚。）
\end{enumerate}

\begin{enumerate}
\def\labelenumi{\arabic{enumi})}
\setcounter{enumi}{12}
\tightlist
\item
    设 D 和 D' 是 \(R^{3}\) 中两条不同的直线，令 \(U = R^{3} - (D \cup D')\)。
\end{enumerate}

\begin{enumerate}
\def\labelenumi{\alph{enumi})}
\item
  证明 U 是道路连通的。
\item
  假设 D 和 D' 不相交；证明对每个点 \(a \in U\)，群 \(\pi_{1}(U, a)\) 是以两个生成元为基的自由群。
\item
  假设 D 和 D' 有公共点；证明对每个点 \(a \in U\)，群 \(\pi_{1}(U, a)\) 是以三个生成元为基的自由群。
\end{enumerate}

\begin{enumerate}
\def\labelenumi{\arabic{enumi})}
\setcounter{enumi}{13}
\tightlist
\item
    设 P 是 \(R^{3}\) 中的一个平面，C 是包含于 P 的圆，D 是 \(R^{3}\) 中的一条直线。令 \(U = R^{3} - (C \cup D)\)，并令 a 为 U 中的一点。
\end{enumerate}

\begin{enumerate}
\def\labelenumi{\alph{enumi})}
\item
  证明 U 是道路连通的。
\item
  假设 D 不与 C 的凸包相交。证明 \(\pi_{1}(U,a)\) 是以两个生成元为基的自由群。
\item
  假设 D 不与 C 相交，但与 C 的凸包相交。证明 \(\pi_{1}(U,a)\) 同构于 \(Z^{2}\)。
\item
  假设 D 与 C 恰好交于一点。证明 \(\pi_{1}(U,a)\) 是以两个生成元为基的自由群。
\item
  假设 D 与 C 交于两点。证明 \(\pi_{1}(U,a)\) 是以三个生成元为基的自由群。
\end{enumerate}

\begin{enumerate}
\def\labelenumi{\arabic{enumi})}
\setcounter{enumi}{14}
\tightlist
\item
  若拓扑空间 X 是连通的，并且对于 X 的每一对不交闭子集 (A, B)，只要 \(X-A\) 和 \(X-B\) 连通，就有 \(X-(A \cup B)\) 连通，则称 X 具有弗拉格门–布劳威尔性质。
\end{enumerate}

\begin{enumerate}
\def\labelenumi{\alph{enumi})}
\item
  设 X 是具有弗拉格门–布劳威尔性质的局部连通拓扑空间。令 (A, B) 是 X 的一对不交闭子集，令 \((a, b)\) 是 \(X-(A \cup B)\) 中的一对点，它们属于 \(X-A\)（分别属于 \(X-B\)）的同一连通分支。证明 a、b 属于 \(X-(A \cup B)\) 的同一连通分支。
\item
  设 X 是连通且局部道路连通的拓扑空间，但不具有弗拉格门–布劳威尔性质。证明，对每个点 \(x \in X\)，存在从 \(\pi_{1}(X, x)\) 到 Z 的满射同态。
\item
  设 X 是连通、局部道路连通且道路单连通的拓扑空间。证明 X 具有弗拉格门–布劳威尔性质。
\end{enumerate}

d) 设 X 是分离且局部连通的拓扑空间，满足对每个点 \(x \in \mathrm{X}\)，空间 \(\mathrm{X} - \{x\}\) 具有弗拉格门–布劳威尔性质。令 A 是同胚于 \([0,1]\) 的 X 的子集；证明 \(\mathrm{X} - \mathrm{A}\) 连通。（令 \(c: \mathbf{I} \to \mathrm{A}\) 为同胚。反设并构造一个序列 \((\mathrm{J}_n)\)，使得对每个 \(n\)，\(\mathrm{J}_n\) 是 \(\mathbf{I}\) 中长度为 \(2^{-n}\) 的区间，包含于 \(\mathrm{J}_{n-1}\) 中（若 \(n \geqslant 1\)），并且 \(\mathrm{X} - c(\mathrm{J}_n)\) 不连通；令 \(x\) 为 X 中属于所有集合 \(c(\mathrm{J}_n)\) 的唯一点。然后证明 \(\mathrm{X} - \{x\}\) 不连通。）

\begin{enumerate}
\def\labelenumi{\alph{enumi})}
\setcounter{enumi}{4}
\tightlist
\item
  设 \(n\) 是整数且 \(\geqslant 1\)，令 A 是 \(\mathbf{S}_n\) 中同胚于 \([0,1]\) 的子集。证明 \(\mathbf{S}_n - \mathbf{A}\) 连通。
\end{enumerate}

\begin{enumerate}
\def\labelenumi{\arabic{enumi})}
\setcounter{enumi}{15}
\tightlist
\item
  设 A 是 \(S_{2}\) 中同胚于 \(S_{1}\) 的子集。令 a、b 是 A 的不同点；记 \(A'\) 和 \(A''\) 为 \(A-\{a,b\}\) 的连通分支，并置 \(X=S_{2}-\{a,b\}\)、\(U=X-A'\) 和 \(V=X-A''\)。
\end{enumerate}

\begin{enumerate}
\def\labelenumi{\alph{enumi})}
\item
  证明 U 和 V 连通。（使用练习 15。）
\item
  证明对每个点 \(x \in U\)，由 U 包含于 X 所诱导的从 \(\pi_{1}(U, x)\) 到 \(\pi_{1}(X, x)\) 的同态是平凡的。
\item
  证明 \(S_{2}-A\) 恰有两个连通分支，且它们的边界等于 A。
\item
  设 S 是 \(R^{2}\) 中同胚于 \(S_{1}\) 的子集；证明 \(R^{2}-S\) 恰有两个连通分支，且它们的边界等于 S（若尔当定理）。
\end{enumerate}

\begin{enumerate}
\def\labelenumi{\arabic{enumi})}
\setcounter{enumi}{16}
\tightlist
\item
  我们将 \(S_{1}\) 识别为模长为 1 的复数集，并将球面 \(S_{3}\) 识别为 \(C^{2}\) 的子空间，由点 \((z,w)\) 组成，这些点满足 \(|z|^{2} + |w|^{2} = 1\)。令 B、C 分别为 \(S_{3}\) 中由 \(|z| \leqslant |w|\) 和 \(|z| \geqslant |w|\) 定义的子空间。
\end{enumerate}

\begin{enumerate}
\def\labelenumi{\alph{enumi})}
\item
  证明 B 和 C 同胚于 \(B_{2} \times S_{1}\)。
\item
  对于每一对互素整数 \((m,n)\)，记 \(K_{m,n}\) 为从 \(S_{1}\) 到 \(S_{3}\) 的映射 \(u\mapsto(u^{m},u^{n})\) 的像（“环面纽结”）。
\item
  证明 \(S_{3}-K_{m,n}\) 连通，且其基本群具有表现 \(\langle x,y;x^{m}=y^{n}\rangle\)。
\end{enumerate}

d) 证明，存在同胚 \(\varphi\)，它从 \(\mathbf{S}_3\) 到自身，并且满足 \(\varphi(\mathrm{K}_{m,n}) = \mathrm{K}_{1,0}\)，当且仅当 \(|m| \leqslant 1\) 或 \(|n| \leqslant 1\)。

\begin{enumerate}
\def\labelenumi{\arabic{enumi})}
\setcounter{enumi}{17}
\tightlist
\item
  设 \((m,n)\) 是满足 1 \textless{} m \textless{} n 的一对自然整数。令 G 为由表现 \(\langle x,y;x^{m}=y^{n}\rangle\) 定义的群。令 \(z=x^{m}\)，令 C 为 G 中由 z 生成的子群。
\end{enumerate}

\begin{enumerate}
\def\labelenumi{\alph{enumi})}
\item
  证明 C 包含于 G 的中心。
\item
  证明 G/C 同构于自由积 \((\mathbf{Z}/m\mathbf{Z})*(\mathbf{Z}/n\mathbf{Z})\)。
\item
  推出 C 是 G 的中心。
\item
  设 \((p,q)\) 是满足 1 \textless{} p \textless{} q 的一对自然整数。假设群 \(G'\) 由表现 \(\langle x,y; x^{p} = y^{q}\rangle\) 定义且同构于 G。证明 \((p,q) = (m,n)\)。
\item
  证明 G 对其导出群的商同构于 Z。
\end{enumerate}

\begin{enumerate}
\def\labelenumi{\arabic{enumi})}
\setcounter{enumi}{18}
\tightlist
\item
  设 G 是图；记 S 为其顶点集，A 为其有向边集，\(|\mathbf{G}|\) 为其几何实现。记 \(j\colon \mathrm{S}\to |\mathrm{G}|\) 和 \(p\colon \mathbf{I}\times \mathbf{A}\to |\mathbf{G}|\) 为典范映射。假设集合 S 和 A 都是有限的。
\end{enumerate}

一个映射 \(f\colon|G|\to\mathbf{R}^{n}\) 称为仿射的（分别为分段仿射的），若对每条边 \(a\in A\)，映射 \(t\mapsto f\circ p(a,t)\) 是仿射的（分别为分段仿射的）。（参见 III，第 331 页，练习 6。）

\begin{enumerate}
\def\labelenumi{\alph{enumi})}
\item
   证明，存在从 \(|G|\) 到 \(R^{3}\) 的单射分段仿射映射。（先证明存在整数 \(n \geqslant 1\) 以及从 \(|G|\) 到 \(R^{n}\) 的单射分段仿射映射。）
\item
   对每个连续单射映射 \(f\colon|G|\to R^{n}\)，证明存在连续单射映射 \(g\colon|G|\to R^{n}\)，它是分段仿射的并且与 f 同伦。
\item
   假设映射 \((o,t)\colon\mathrm{A}\to\mathrm{S}\times\mathrm{S}\) 是单射。令 f 为从 S 到 R 的单射映射。证明，存在唯一的从 \(|G|\) 到 \(R^{3}\) 的仿射映射 F，它把每个顶点 \(s\in S\) 映到点 \((f(s),f(s)^{2},f(s)^{3})\)（该点属于 \(R^{3}\)）。证明 F 是单射。
\end{enumerate}

我们称图 G 是平面的，如果存在从 \(|G|\) 到 \(R^{2}\) 的连续单射映射。

d) 假设图 \(\mathbf{G}\) 是平面的，且映射 \((o,t)\colon\mathbf{A}\to\mathbf{S}\times\mathbf{S}\) 是单射。令 \(f\colon|\mathbf{G}|\to\mathbf{R}^{2}\) 是连续单射映射。证明，存在单射仿射映射 \(g\colon|\mathbf{G}|\to\mathbf{R}^{2}\)，它与 \(f\) 同伦。

\begin{enumerate}
\def\labelenumi{\arabic{enumi})}
\setcounter{enumi}{19}
\tightlist
\item
   设 \(f \colon \mathbf{I} \to \mathbf{R}^2\) 是一个分段仿射圈，且其在 [0,1[ 上的限制是单射；记 \(C = f(\mathbf{I})\)。
\end{enumerate}

\begin{enumerate}
\def\labelenumi{\alph{enumi})}
\item
   证明 \(\mathbf{R}^2 -\mathbf{C}\) 至多有两个连通分支。
\item
   对于 \(x \in R^{2} - C\) 和 \(v \in S_{1}\)，令 \(n_{v}(x)\) 为 \((\mathbf{R}^{2} - \mathbf{C}) \cap (x + \mathbf{R}_{+}v)\) 的连通分支数。证明，存在唯一的局部常值映射 \(n: R^{2} - C \to \{0,1\}\)，使得 \(n(x) \equiv n_{v}(x) \pmod{2}\) 对所有 \(v \in S_{1}\) 都成立。
\item
   推出 \(R^{2}-C\) 恰有两个连通分支，且它们的边界等于 C。
\item
   设 \(g: I \to R^{2}\) 是单射分段仿射映射，满足 \(\overline{g}^{-1}(C) = \{0,1\}\)；置 \(P = g(I)\)，\(D = C \cup P\)。证明 \(C - (C \cap P)\) 有两个连通分支；记它们为 \(P_{1}\) 和 \(P_{2}\)。证明 \(R^{2} - D\) 有三个连通分支，且它们的边界分别等于 C、\(P_{1} \cup P\) 和 \(P_{2} \cup P\)。
\item
   设 G 是有限平面图；令 S 为其顶点集，A 为其定向边集。令 \(f\colon|G|\to\mathbf{R}^{2}\) 是连续单射分段仿射映射，令 P 为其像。证明 \(\mathbf{R}^{2}-P\) 的连通分支数等于 \(1+\mathrm{Card(S)}-\mathrm{Card(A)}+\mathrm{Card}(\pi_{0}(\mathrm{G}))\)。
\end{enumerate}

\begin{enumerate}
\def\labelenumi{\arabic{enumi})}
\setcounter{enumi}{20}
\item
  令 K 为与如下箭图相关的图：其顶点集为 \(\{1,2,3\} \times \{0,1\}\)，箭集为形如 \(((a,0),(b,1))\) 的所有对，其中 \(a,b\in \{1,2,3\}\)，而其源映射和靶映射分别由第一、第二投影导出（“库拉托夫斯基图”）。证明图 K 不是平面的。
\item
  设 \(n\) 是自然数，令 \(\mathrm{K}_n\) 为顶点集基数为 \(n\) 的完全图（II，第 220 页，练习 5）。
\end{enumerate}

\begin{enumerate}
\def\labelenumi{\alph{enumi})}
\item
  证明图 \(K_{n}\) 在 \(1 \leqslant n \leqslant 4\) 时是平面的。
\item
  证明图 \(K_{5}\) 不是平面的。
\end{enumerate}

\begin{enumerate}
\def\labelenumi{\arabic{enumi})}
\setcounter{enumi}{22}
\tightlist
\item
   设 \(u\)、\(v\) 是 \(\mathbf{S}_1\) 中满足 \(\mathcal{I}(v / u) > 0\) 的两个点；定义族 \((a_j)_{j \in \mathbf{Z} / 4\mathbf{Z}}\)，其中 \(a_1 = u\)、\(a_2 = v\)、\(a_3 = -u\)、\(a_4 = -v\)。记 \(\mathrm{Q}_1\)、\(\mathrm{Q}_2\) 为线段 \([a_1, a_3]\) 和 \([a_2, a_4]\)，它们位于单位圆盘 \(\mathbf{B}_2\) 中，并置 \(\mathrm{Q} = \mathrm{Q}_1 \cup \mathrm{Q}_2\)。
\end{enumerate}

\begin{enumerate}
\def\labelenumi{\alph{enumi})}
\tightlist
\item
  证明，对每个 \(j \in Z/4Z\)，存在唯一的 \(A_{j} \in \pi_{0}(B_{2}-Q)\) 的道路连通分支，其闭包包含 \(a_{j}\) 和 \(a_{j-1}\)。证明族 \((A_{j})\) 的并等于 \(B_{2}-Q\)。
\end{enumerate}

  置 \(C = B_{2} \times [-1, 1]\)、\(N = (0, 0, 1)\)、\(S = (0, 0, -1)\)。令 \(f: [-1, 1] \to R_{+}\) 为连续映射，满足 \(f(-1) = f(1) = 0\)，且 \(0 < f(t) < 1\) 对所有 \(t \in ]-1; 1[\) 都成立。于是记 \(L_{1}\)、\(L_{2}\) 为从 \([-1, 1]\) 到 C 的映射 \(t \mapsto (ta_{1}, 0)\) 和 \(t \mapsto (ta_{2}, f(t))\) 的像；最后置 \(L = L_{1} \cup L_{2}\)。

\begin{enumerate}
\def\labelenumi{\alph{enumi})}
\setcounter{enumi}{1}
\item
  证明 C-L 道路连通。
\item
  若一条道路 \(c\) 位于 \(\mathbf{C} - \mathbf{L}\) 中，起点为 \(\mathbf{N}\) 且终点为 \(\mathbf{S}\)，并且存在相邻道路 \(c_1\)、\(c_2\) 位于 \(\mathbf{B}_2\) 中，使得 \(c\) 是由路径 \(t \mapsto (c_1(t), 1)\)、\(t \mapsto (c_1(1), 1 - 2t)\) 和 \(t \mapsto (c_2(t), -1)\) 拼接而成的，则称 c 为直的。证明此时 \(c_1(1) \not\in \mathbb{Q}\)；记 \(\mathrm{A}_c\) 为包含该点的 \(\mathbf{B}_2 - \mathbf{Q}\) 的道路连通分支。
\end{enumerate}

d) 设 \(c\) 和 \(c'\) 是 \(\mathrm{C} - \mathrm{L}\) 中起点为 \(\mathrm{N}\)、终点为 \(\mathrm{S}\) 的道路；假设它们是直的。证明它们严格同伦，当且仅当 \(\mathrm{A}_c = \mathrm{A}_{c'}\)。

\begin{enumerate}
\def\labelenumi{\alph{enumi})}
\setcounter{enumi}{4}
\tightlist
\item
  对每个 \(j \in \mathbf{Z}/4\mathbf{Z}\)，令 \(c_j\) 是 \(C-L\) 中起点为 \(N\)、终点为 \(S\) 的直道路，且满足 \(A_{c_j} = A_j\)。证明关系式
\end{enumerate}

\[
[ c _ {1} ] [ c _ {4} ] ^ {- 1} = [ c _ {2} ] [ c _ {3} ] ^ {- 1}
\]

在群 \(\pi_1(\mathrm{C} - \mathrm{L},\mathrm{N})\)

\begin{enumerate}
\def\labelenumi{\alph{enumi})}
\setcounter{enumi}{5}
\tightlist
\item
   证明，从以 \(x, y\) 为两个生成元的自由群到 \(\pi_1(\mathrm{C} - \mathrm{L}, \mathrm{N})\) 的典范同态，将 \(x\) 和 \(y\) 分别映到 \([c_1][c_2]^{-1}\) 和 \([c_2][c_3]^{-1}\)，是群同构。（将范坎彭定理应用于 \(\mathrm{C} - \mathrm{L}\) 的覆盖，该覆盖由 \(\mathbf{R}^3\) 中包含集合 \(\{a_1, a_2, \mathrm{N}\}\) 中两点的三个向量平面所界定的八个闭集组成。）
\end{enumerate}

\begin{enumerate}
\def\labelenumi{\arabic{enumi})}
\setcounter{enumi}{23}
\tightlist
\item
   设 m 是 \(\geqslant 1\) 的整数，令 \((\theta_{j})_{1\leqslant j\leqslant m}\) 是实数序列，满足 \(0\leqslant\theta_{1}<\theta_{2}<\cdots<\theta_{n}<2\pi\)。对 \(j\in Z/mZ\)，置 \(a_{j}=\exp(i\theta_{k})\)，其中 k 是属于类 j 的 \(\{1,\ldots,n\}\) 中唯一元素；再记 \(Q_{j}\) 为线段 \([0,a_{j}]\)，它位于 \(B_{2}\) 中。最后置 \(Q=Q_{1}\cup\cdots\cup Q_{m}\)。
\end{enumerate}

\begin{enumerate}
\def\labelenumi{\alph{enumi})}
\item
  对 \(j \in Z/mZ\)，证明存在唯一的 \(A_{j}\)，它是 \(B_{2}-Q\) 的道路连通分支，并包含所有 \(b \in S_{1}\)，其中满足 \(\mathcal{I}(b/a_{j}) > 0\) 且 \(\mathcal{I}(b/a_{j+1}) < 0\)。证明族 \((A_{j})\) 的成员两两不交，且其并等于 \(B_{2}-Q\)。
\item
   置 \(C = B_{2} \times [-1, 1]\)、\(N = (0, 0, 1)\)、\(S = (0, 0, -1)\)，\(L = Q \times \{0\}\)。若 \(C-L\) 中一条起点为 N、终点为 S 的道路 c 满足存在相邻的道路 \(c_{1}\)、\(c_{2}\) 位于 \(B_{2}\) 中，使得 c 是由路径 \(t \mapsto (c_{1}(t), 1)\)、\(t \mapsto (c_{1}(1), 1 - 2t)\) 和 \(t \mapsto (c_{2}(t), -1)\) 拼接而成的，则称其为直的。证明此时 \(c_{1}(1) \notin Q\)；记 \(A_{c}\) 为包含该点的 \(B_{2}-Q\) 的道路连通分支。
\item
   令 c 和 \(c'\) 是 \(C-L\) 中起点为 N、终点为 S 的道路；假设它们是直的。证明它们严格同伦，当且仅当 \(A_{c} = A_{c'}\)。
\end{enumerate}

d) 对每个 \(j \in \mathbf{Z}/m\mathbf{Z}\)，令 \(c_j\) 是 \(C-L\) 中起点为 N、终点为 S 的直圈，且满足 \(\mathrm{A}_{c_j} = \mathrm{A}_j\)。证明从自由群 \(\mathrm{F}(\mathbf{Z}/m\mathbf{Z})\) 到 \(\pi_1(\mathrm{C}-\mathrm{L},\mathrm{N})\) 的同态将 \(x_j\) 映为 \([c_j][c_{j+1}]^{-1}\)，对每个 \(j \in \mathbf{Z}/m\mathbf{Z}\)，它是满射，且其核是包含乘积 \(x_1x_2 \ldots x_m\) 的最小正规子群。

\begin{enumerate}
\def\labelenumi{\arabic{enumi})}
\setcounter{enumi}{24}
\tightlist
\item
   置 \(C = B_{2} \times [-1, 1]\)。令 K 是 \(R^{3}\) 的闭子空间，令 P 为它在映射 \((x, y, z) \mapsto (x, y)\) 下的像。假设存在有限集 I、点族 \((p_{i})_{i \in \mathrm{I}}\)（点属于 \(R^{2}\)）、严格正实数族 \((r_{i})_{i \in \mathrm{I}}\) 以及自然数族 \((m_{i})_{i \in \mathrm{I}}\)，满足以下性质，其中 \(\varphi_{i}\) 表示从 \(R^{3}\) 到 \(R^{3}\) 的映射 \(x \mapsto (p_{i}, 0) + r_{i}x\)，并且 \(C_{i} = \varphi_{i}(C) \)：
\end{enumerate}

- 各集合 \(C_{i}\) 两两不交；

- 对每个 \(i \in \mathrm{I}\)，集合 \(\mathrm{L}_i = \overline{\varphi}_i^{-1}(\mathrm{C}_i \cap \mathrm{K})\) 是练习 23 中定义的子空间 \(\mathrm{L}\)，若 \(m_i = 0\)，以及练习 24（其中置 \(m = m_i\)），若 \(m_i \geqslant 1\)。

- 集合 \(\mathrm{K}' = \mathrm{K} - \bigcup_{i} (\mathring{\mathrm{C}}_i \cap \mathrm{K})\) 是有限族闭线段的并，且这些线段在平面 \(\mathbf{R}^2 \times \{0\}\) 中两两不交。

令 \(\mathrm{N} = (x_{\mathrm{N}}, y_{\mathrm{N}}, z_{\mathrm{N}})\) 和 \(\mathrm{S} = (x_{\mathrm{S}}, y_{\mathrm{S}}, z_{\mathrm{S}})\) 为 \(R^{3}\) 中满足 \(z_{N} > \sup(r_{i})_{i \in I}\) 且 \(z_{\mathrm{S}} < -\sup(r_{i})_{i \in I}\) 的点。

\begin{enumerate}
\def\labelenumi{\alph{enumi})}
\tightlist
\item
  称道路 \(c\) 在 \(\mathbf{R}^3-K\) 中且起点为 N、终点为 S 为直的，若存在相邻的道路 \(c_1\) 和 \(c_2\) 位于 \(\mathbf{R}^2\) 中，使得 \(c\) 是以下道路的拼接：\(t \mapsto (c_1(t), z_{\mathrm{N}})\)、\(t \mapsto (c_1(1), (1 - t)z_{\mathrm{N}} + tz_{\mathrm{S}})\) 以及 \(t \mapsto (c_2(t), z_{\mathrm{S}})\)。
\end{enumerate}

证明此时点 \(c_1(1)\) 不属于 P；记 \(A_c\) 为包含该点的 \(\mathbf{R}^2 - \mathbf{P}\) 的道路连通分支。证明两条直道路 \(c\) 和 \(c'\) 严格同伦，当且仅当 \(A_c = A_{c'}\)。

   对于 \(\alpha\) 这个 \(\mathbf{R}^2-\mathbf{P}\) 的道路连通分支，我们选取一条道路 \(c(\alpha)\)，其位于 \(\mathbf{R}^2-\mathbf{K}\) 中，起点为 N、终点为 S，且为直道路，使得 \(\mathrm{A}_{c(\alpha)} = \alpha\)。

\begin{enumerate}
\def\labelenumi{\alph{enumi})}
\setcounter{enumi}{1}
\item
   设 \(i \in I\) 满足 \(m_{i} = 0\)，且 \(\overline{\varphi}_{i}^{-1}(K \cap C_{i})\) 是练习 23 中定义的集合 L；有 \(\overline{\varphi}_{i}^{-1}(P \times \{0\} \cap C_{i}) = Q \times \{0\}\)。对所有 \(j \in Z/4Z\)，令 \(\alpha_{i,j}\) 为 \(R^{2} - P\) 的道路连通分支，它包含 \(\varphi_{i}\) 所映的连通分支 \(A_{j}\)，该分支属于 \(B_{2} - P\)。证明关系式 \([c(\alpha_{i,1})][c(\alpha_{i,4})]^{-1} = [c(\alpha_{i,2})][c(\alpha_{i,3})]^{-1}\) 在 \(\pi_{1}(R^{3} - K, N)\) 中成立。
\item
   令 \(\omega\) 是 \(\pi_0(\mathbf{R}^2-\mathbf{P})\) 中的元素。令 \(\lambda_\omega\colon \mathrm{F}(\pi_0(\mathbf{R}^2-\mathbf{P}))\to \pi_1(\mathbf{R}^3-\mathbf{K},\mathbf{N})\) 为唯一的群同态，对于每个 \(\alpha \in \pi_0(\mathbf{R}^2-\mathbf{P})\)，它把元素 \(x_\alpha\) 映为 \([c(\alpha)][c(\omega)]^{-1}\)。
\end{enumerate}

证明同态 \(\lambda_{\omega}\) 是满射，其核是 \(\mathrm{F}(\pi_0(\mathbf{R}^2 - \mathbf{P}))\) 的包含元素 \(x_{\omega}\) 以及 \(x_{\alpha_{i,1}}^{-1}x_{\alpha_{i,2}}x_{\alpha_{i,3}}^{-1}x_{\alpha_{i,4}}\)（其中 \(i \in \mathrm{I}\) 满足 \(m_i = 0\)）的最小正规子群。

\begin{enumerate}
\def\labelenumi{\alph{enumi})}
\setcounter{enumi}{3}
\item
   对于每条线段 \(s\)，它属于 \(\pi_{0}(\mathrm{K}^{\prime})\)，选取元素 \(\alpha_{s}\) 和 \(\alpha_{s}^{\prime}\)，它们属于 \(\pi_{0}(\mathbf{R}^{2}-\mathbf{P})\)，使得 \(R^{2}-P\) 中闭包包含 \(s\) 的连通分支集合等于 \(\{\alpha_{s},\alpha_{s}^{\prime}\}\)。令 \(\mu\colon\mathrm{F}(\pi_{0}(\mathrm{K}^{\prime}))\to\pi_{1}(\mathbf{R}^{3}-\mathbf{K},\mathbf{N})\) 为唯一的群同态，对于 \(s\in\pi_{0}(\mathrm{K}^{\prime})\)，它把元素 \(x_{s}\) 映为 \([c(\alpha_{s})][c(\alpha_{s}^{\prime})]^{-1}\)。证明 \(\mu\) 是满射。对于 \(i\in I\)，其中 \(m_{i}=0\)，以及 \(j\in\{1,2,3,4\}\)，令 \(s_{i,j}\) 为 \(\pi_{0}(\mathrm{K}^{\prime})\) 中满足 \(\varphi_{i}(a_{j})\) 与 \(s_{i,j}\) 相交的唯一线段。证明存在元素 \(u_{i},v_{i},w_{i}\)，它们属于 \(\{-1,1\}\)，使得 \(\mu(x_{s_{i,4}})=\mu(x_{s_{i,2}}^{w_{i}})\) 且 \(\mu(x_{s_{i,3}})=\mu(x_{s_{i,2}}^{-u_{i}}x_{s_{i,1}}^{v_{i}}x_{s_{i,2}}^{u_{i}})\)。证明 \(\mu\) 的核是 \(\mathrm{F}(\pi_{0}(\mathrm{K}^{\prime}))\) 的包含元素 \(x_{s_{i,3}}x_{s_{i,2}}^{-u_{i}}x_{s_{i,1}}^{-v_{i}}x_{s_{i,2}}^{u_{i}}\) 和 \(x_{s_{i,4}}x_{s_{i,2}}^{-w_{i}}\)（对 \(i\in I\)）的最小正规子群。
\item
  Set
\end{enumerate}

\[
k = \operatorname{Card} (\pi_ {0} (\mathrm{K})) - \operatorname{Card} (\{i \in \mathrm{I} \mid m _ {i} > 0 \}) + \frac {1}{2} \sum_ {i \in \mathrm{I}} m _ {i}.
\]

   证明 \(k \in N\)。证明 \(\pi_{1}(\mathbf{R}^{3} - \mathbf{K})\) 对其导出群的商同构于 \(Z^{k}\)。

\begin{enumerate}
\def\labelenumi{\arabic{enumi})}
\setcounter{enumi}{25}
\tightlist
\item
  在数值平面 \(\mathbf{R}^2\) 中，令 \(\mathrm{D}_0\) 为中心为 \((0,0)\)、半径为 1 的圆盘，令 \(\mathrm{C}_0\) 为其边界；置 \(a_0 = (-1,0)\)。令 \(g\) 是满足 \(\geqslant 1\) 的自然整数，令 \(r\) 是严格正实数，令 \((v_1,\ldots,v_g)\) 是实数序列，满足 \(r < \inf (v_{1} + 1,(v_{2} - v_{1}) / 2,\dots,(v_{g} - v_{g - 1}) / 2,1 - v_{g})\)。对于所有 \(j\in \{1,\ldots,g\}\)，记 \(\mathrm{D}_j\) 为中心为 \((0,v_j)\) 的闭圆盘
\end{enumerate}

半径为 \(r\)，并令 \(\mathrm{C}_j\) 为中心为 \((0,v_j)\)、半径为 \(r\) 的圆；再置 \(a_{j}=(-r,v_{j})\)。令 X 为空间 \(\mathrm{D}=\bigcup_{j=1}^{g}\mathring{\mathrm{D}}_{j}\)。

  对每个 \(j \in \{1, \ldots, j\}\)，令 \(u_j\) 为 \(\varpi(\mathrm{X})\) 中一条起点为 \(a_0\)、终点为 \(a_j\) 且像为线段 \([a_0, a_j]\) 的道路的道路类；也记 \(c_j \in \pi_1(\mathrm{X}, a_j)\) 为由 \(t \mapsto (-r \cos(2\pi t), -r \sin(2\pi t) + v_j)\) 给出的圈的道路类，并令 \(\gamma_j = u_j c_j u_j^{-1}\)。

\begin{enumerate}
\def\labelenumi{\alph{enumi})}
\item
  令 \(\gamma_{0}\) 为 X 中由 \(t \mapsto (-\cos(2\pi t), -\sin(2\pi t))\) 给出的圈的道路类。证明有 \(\gamma_{0} = \gamma_{1} \ldots \gamma_{g}\)。
\item
  证明从以 \(x_{1},\ldots,x_{j}\) 为生成元的自由群到 \(\pi_{1}(X,a_{0})\) 的唯一同态（将 \(x_{j}\) 映为 \(\gamma_{j}\)）是群同构。（将范坎彭定理应用于 X 与半平面 \(R_{+}\times R\) 和 \(R_{-}\times R\) 的交所定义的覆盖。）
\item
  令 \(f \colon \mathrm{X} \to \mathbf{R}_+\) 是连续函数，满足 \(f^{-1}(0) = \bigcup_{j=0}^{g} \mathrm{C}_j\)，令 Y 为 \(\mathbf{R}^3\) 中的点 \((x, y, z)\) 所成的子空间，这些点满足 \(z^2 = f(x, y)\)。记 \(i_+\)（分别为 \(i_-\)）为从 X 到 Y 的连续映射，分别由 \((x, y) \mapsto (x, y, \sqrt{f(x, y)})\)（分别由 \((x, y) \mapsto (x, y, -\sqrt{f(x, y)})\)）给出。对于 X 中的每个道路类 \(\gamma\)，在 \(\varpi(\mathrm{Y})\) 中定义道路类 \(\gamma^+ = (i_+)_{*}(\gamma)\) 和 \(\gamma^- = (i_-)_*(\gamma)\)。
\end{enumerate}

对于每个 \(j \in \{1, \ldots, g\}\)，最后在 \(\pi_1(Y, a_0)\) 中定义元素 \(\delta_j = u_j^+ u_j^-\)。证明有关系 \(\gamma_j^+ \delta_j = \delta_j \gamma_j^-\)（\(1 \leqslant j \leqslant g\)），以及 \(\gamma_1^+ \ldots \gamma_g^+ = \gamma_1^- \ldots \gamma_g^-\)。

\begin{enumerate}
\def\labelenumi{\alph{enumi})}
\setcounter{enumi}{3}
\tightlist
\item
  令 \(\varphi\) 为从具有 2g 个生成元 \(x_{1},\ldots ,x_{g},y_{1},\ldots ,y_{g}\) 的自由群到 \(\pi_1(\mathrm{Y},a_0)\) 的唯一同态，它把 \(x_{j}\) 映为 \(\gamma_j^+\)，把 \(y_{j}\) 映为 \(\delta_{j}\)。证明同态 \(\varphi\) 是满射，且其核是包含元素
\end{enumerate}

\[
(x _ {1} ^ {- 1} y _ {1} x _ {1}) (x _ {2} ^ {- 1} y _ {2} x _ {2}) \dots (x _ {g} ^ {- 1} y _ {g} x _ {g}) (y _ {1} \dots y _ {g}) ^ {- 1}.
\]

（将范坎彭定理应用于 Y 与半空间 \(R^{2} \times R_{+}\) 和 \(R^{2} \times R_{-}\) 的交所定义的覆盖。）

\begin{enumerate}
\def\labelenumi{\alph{enumi})}
\setcounter{enumi}{4}
\tightlist
\item
  令 \(\varphi'\) 为从具有 2g 个生成元 \(x_1, \ldots, x_g, y_1, \ldots, y_g\) 的自由群到 \(\pi_1(Y, a_0)\) 的唯一同态，对于每个 \(j\)，它把 \(x_j\) 映为
\end{enumerate}

\[
(\gamma_ {1} ^ {+} \dots \gamma_ {j - 1} ^ {+}) ^ {- 1} \delta_ {j} (\gamma_ {1} ^ {+} \dots \gamma_ {j - 1} ^ {+})
\]

并把 \(y_{j}\) 映为 \(\gamma_{j}^{+}\)。证明同态 \(\varphi'\) 是满射，且其核是包含元素

\[
\left(x _ {1} ^ {- 1} y _ {1} ^ {- 1} x _ {1} y _ {1}\right) \dots \left(x _ {g} ^ {- 1} y _ {g} ^ {- 1} x _ {g} y _ {g}\right).
\]

\begin{enumerate}
\def\labelenumi{\arabic{enumi})}
\setcounter{enumi}{26}
\tightlist
\item
  令 D 为 C 中的单位圆盘，令 X 为 D 的一个子空间，且是 \(S_{1}\) 的一个邻域。令 R 为 X 上的最粗等价关系：对于 \(u, v \in S_{1}\)，当且仅当点 \(f(u)\) 和点 \(f(v)\) 满足 \(f(u) = f(v)\) 时，识别它们。令 \(p\colon X \to X / R\) 为典范满射。令 \(T\) 为紧拓扑空间，令 \(f\colon S_1 \to T\) 为连续满射。置 \(a = f(1)\)，并记 \(\alpha \in \pi_1(T, a)\) 为由 \(t \mapsto f(e^{2\pi it})\) 给出的圈的道路类，它位于 \(T\) 中。
\end{enumerate}

\begin{enumerate}
\def\labelenumi{\alph{enumi})}
\item
  证明，存在唯一的从 T 到 X/R 的映射 \(\sigma\)，使得 \(\sigma \circ f = p|_{\mathbf{S}_1}\)。证明 \(\sigma\) 将 T 同胚到其像 \(p(\mathbf{S}_1)\)。
\item
  令 r 是满足 0 \textless{} r \textless{} 1 的实数；假设 X 是由 \(z \in D\) 且满足 \(r \leqslant |z| \leqslant 1\) 的点所成的集合。证明 X/R 同胚于 f 的映射柱。推出同态 \(\sigma_{*}\) 是从 \(\pi_{1}(T, a)\) 到 \(\pi_{1}(X/R, p(1))\) 的群同构。
\item
  假设 X = D。证明 X/R 同胚于 f 的映射锥。推出群同态 \(\sigma_{*}\colon\pi_{1}(\mathrm{T},a)\to\pi_{1}(\mathrm{X}/\mathrm{R},p(1))\) 是满射，其核是 \(\pi_{1}(\mathrm{T},a)\) 中包含 \(\alpha\) 的最小正规子群。
\item
  令 r 是满足 0 \textless{} r \textless{} 1 的实数；令 \(X_{1}\) 为由 \(z \in X\) 且满足 \(|z| \leqslant r\) 的 z 所成的集合，令 \(X_{2}\) 为由 \(z \in D\) 且满足 \(r \leqslant |z| \leqslant 1\) 的 z 所成的集合；假设 \(X_{2} \subset X\)。令 \(\beta \in \pi_{1}(X_{2}, r)\) 为由 \(t \mapsto re^{2\pi it}\) 给出的圈的道路类；令 \(\delta \in \varpi_{p(r), p(1)}\) 为由 \(t \mapsto p((1 - t)r + t)\) 给出的道路的道路类。
\end{enumerate}

  证明，存在唯一的群同态 \(\mu\)，从 \(\pi_1(\mathrm{T},a)*\pi_1(\mathrm{X}_2,r)\) 到 \(\pi_1(\mathrm{X} / \mathrm{R},p(1))\)，它与同态 \(\sigma_*\) 在 \(\pi_1(\mathrm{T},a)\) 上一致，并与同态 \(v\mapsto \delta^{-1}p_{*}(v)\delta\) 在 \(\pi_1(\mathrm{X}_2,p(r))\) 上一致。证明 \(\mu\) 是满射，且其核是包含 \(\alpha \beta^{-1}\) 的最小正规子群。

\begin{enumerate}
\def\labelenumi{\arabic{enumi})}
\setcounter{enumi}{27}
\tightlist
\item
  设同伦 \(\sigma: \mathrm{X} \times \mathbf{I} \to \mathrm{X}\) 为同痕，若对每个 \(t \in \mathbf{I}\)，映射 \(x \mapsto \sigma(x, t)\) 是由 \(\sigma\) 诱导的从 X 到自身的同胚。
\end{enumerate}

\begin{enumerate}
\def\labelenumi{\alph{enumi})}
\item
  令 m 是自然整数，令 \(f: X \to R^{n}\) 为连续映射。证明，存在一个同痕，其起点是 \(X \times R^{n}\) 的恒等映射，终点是从 \(X \times R^{n}\) 到 \(X \times R^{n}\) 的映射 \((x, y) \mapsto (x, y - f(x))\)。
\item
  令 A、B 是 X 的子空间；若存在同痕 \(\sigma\colon\mathbf{X}\times\mathbf{I}\to\mathbf{X}\)，其起点是 X 上的恒等映射，且满足 \(f(\mathbf{A}\times\{1\})=\mathbf{B}\)，则称 A 与 B 同痕。
\end{enumerate}

证明关系“ A 与 B 同痕”是在 X 的子空间集合（分别在 X 的闭子空间集合）上的等价关系。

\begin{enumerate}
\def\labelenumi{\alph{enumi})}
\setcounter{enumi}{2}
\tightlist
\item
  设 m、n 是自然数，A 是 \(R^{m}\) 的闭子空间，B 是 \(R^{n}\) 的闭子空间。假设 A 与 B 同胚。证明 \(A \times \{0\}\) 和 \(\{0\} \times B\) 是 \(R^{m+n}\) 的同痕子空间。
\end{enumerate}

\begin{enumerate}
\def\labelenumi{\arabic{enumi})}
\setcounter{enumi}{28}
\tightlist
\item
  设 A 是 \(R^{2}\) 的闭子集，并且同胚于 R 的一个闭子集 B。
\end{enumerate}

\begin{enumerate}
\def\labelenumi{\alph{enumi})}
\item
  证明 \(A \times \{0\}\) 和 \(\{(0,0)\} \times B\) 是 \(R^{3}\) 的同痕子空间。
\item
  假设 \(B \neq R\)；证明 \(R^{2}-A\) 道路连通。
\item
  假设 B = R；证明 \(R^{2}-A\) 恰有两个道路连通分支，且 A 是每个分支的边界。
\item
  设 A 是 \(S_{2}\) 中同胚于 \(S_{1}\) 的子集；证明 \(S_{2} - A\) 恰有两个道路连通分支，且 A 是每个分支的边界。
\item
  设 A 是 \(R^{2}\) 中同胚于 \(S_{1}\) 的子集；证明 \(R^{2}-A\) 恰有两个道路连通分支，且 A 是每个分支的边界。
\end{enumerate}

\begin{enumerate}
\def\labelenumi{\arabic{enumi})}
\setcounter{enumi}{29}
\tightlist
\item
  令 \(v = (v_{1},\ldots ,v_{n})\colon \mathbf{R}^{n}\to \mathbf{R}^{n}\) 是 \(\mathrm{C}^1\) 类映射，令 \(\delta\) 是满足 \(>0\) 的实数；作如下假设：
\end{enumerate}

- \(v\) 的支集包含于 \([- \delta, 1 + \delta] \times \mathring{\mathbf{B}}_{n-1}\) 中；

- 对每个 \(x \in [0, 1 - \delta[\)，都有 \(v_1(x, 0, \ldots, 0) > 0\)，且 \(v_1(1, 0, \ldots, 0) = 0\)；

- 对每个 \(x \in [0,1]\)，都有 \(v_{2}(x,0,\ldots,0) = \cdots = v_{n}(x,0,\ldots,0) = 0\)。

  令 \(\Phi\colon\mathbf{R}^{n}\times\mathbf{R}\to\mathbf{R}^{n}\) 为映射 \((x,t)\mapsto v(x)\) 的积分流（VAR，9.1.3）；对 \(t\in\mathbf{R}\)，令 \(\Phi_{t}\) 为映射 \(x\mapsto\Phi(x,t)\)。

\begin{enumerate}
\def\labelenumi{\alph{enumi})}
\item
  证明存在实数 \(\tau\)，使得 \(\Phi_{\tau}([0,1] \times \{0\}) \subset [1 - \delta, 1] \times \{0\}\)。
\item
  令 \(j\colon[-\delta,1+\delta]\times\mathbf{B}_{n-1}\to\mathbf{R}^{n}\) 是连续单射映射。对每个 \(t\in[0,1]\)，置 \(\mathrm{A}_{t}=j([t,1]\times\{0\})\)。证明对每个 \(t\in]0,1[\)，子空间 \(A_{0}\) 和 \(A_{t}\) 同痕。
\item
  设 P 是 \(\mathbf{R}^{n}\) 中同胚于 \([0,1]\) 且为有限族线段之并的子集。证明 P 与一条线段同痕。
\end{enumerate}

\begin{enumerate}
\def\labelenumi{\arabic{enumi})}
\setcounter{enumi}{30}
\tightlist
\item
  设 \(f = (f_1, \ldots, f_n) \colon \mathbf{I} \to \mathbf{R}^n\) 是 \(\mathbf{C}^1\) 类映射，满足 \(f'(t) \neq 0\) 对所有 \(t \in \mathbf{I}\) 都成立。
\end{enumerate}

\begin{enumerate}
\def\labelenumi{\alph{enumi})}
\item
  假设 \(f_{1}^{\prime}(t) > 0\) 对所有 \(t \in I\) 都成立。对每个 \(j \in \{2, \ldots, n\}\)，证明存在连续映射 \(g_{j} \colon R \to R\)，使得 \(f_{j}(x) = g_{j}(f_{1}(x))\) 对所有 \(x \in I\) 都成立。推出 \(f(\mathbf{I})\) 与一条线段同痕。
\item
  证明存在实数 \(\delta > 0\) 和连续单射映射 \(j\colon [-\delta, 1 + \delta] \times \mathbf{B}_{n-1} \to \mathbf{R}^n\)，使得 \(j(t, 0) = f(t)\) 对所有 \(t \in \mathbf{I} .\) 都成立。由练习 30 推出，对每个 \(\tau \in [0, 1[\)，集合 \(f(\mathbf{I})\) 与 \(f([\tau, 1])\) 同痕。
\item
  证明 \(f(\mathbf{I})\) 与一条线段同痕。
\end{enumerate}

\begin{enumerate}
\def\labelenumi{\arabic{enumi})}
\setcounter{enumi}{31}
\tightlist
\item
  设 A 是 \(\mathbf{R}^3\) 中同胚于 \([0,1]\) 的闭子集。A 的端点是 A 中的两个点 \(p\)，满足 \(\mathrm{A} - \{p\}\) 连通（参见 III，第 280 页，定理 2）；A 的其他点称为内部点。
\end{enumerate}

若 \(p\) 是 A 的一点；称 A 在 \(p\) 处是驯的，若存在 \(p\) 在 \(\mathbf{R}^3\) 中的邻域 V，以及从 V 到 \(\mathbf{R}^3\) 的单位球的同胚，它把 \(\mathrm{A} \cap \mathrm{V}\) 映为一条线段。

\begin{enumerate}
\def\labelenumi{\alph{enumi})}
\item
  设 \((\mathrm{V}_{n})\) 是 \(R^{3}\) 中 p 的邻域的递减序列，满足 \(\bigcap_{n}\mathrm{V}_{n}=\{p\}\)；对每个 n，令 \(a_{n}\) 是 \(V_{n}-A\) 中的一点，并记 \(j_{n}\colon\pi_{1}(\mathrm{V}_{n}-\mathrm{A},a)\to\pi_{1}(\mathrm{V}_{1}-\mathrm{A},a)\) 为由 \(V_{n}\) 包含于 \(V_{1}\) 所诱导的群同态。若 p 是 A 的端点（分别为内部点），证明存在整数 m，使得对每个 \(n\geqslant m\)，同态 \(j_{n}\) 的像归结为单位元（分别为阿贝尔群）。
\item
  假设 A 在每一点处都是驯的。证明存在从 \(R^{3}\) 到自身的同胚，将 A 映为一条线段。推出 \(R^{3}-A\) 连通且道路单连通。
\end{enumerate}

\begin{enumerate}
\def\labelenumi{\arabic{enumi})}
\setcounter{enumi}{32}
\tightlist
\item
  \begin{enumerate}
  \def\labelenumii{(\arabic{enumii})}
  \setcounter{enumii}{5}
  \tightlist
  \item
     设 L 是 \(\mathbf{R}^3\) 的子空间，定义见练习 23，其中取 \(u = (1 - i)/\sqrt{2}\) 和 \(v = (1 + i)/\sqrt{2}\)。令 \(a\)、\(r\) 是满足 \(0 < 2r < a < 1 - 2r\) 的实数。置 \(p_1 = (0, -a)\)、\(p_2 = (0, 0)\)、\(p_3 = (0, a)\)；\(q_1 = (-1, -a)\)、\(q_2 = (-1, 0)\)、\(q_3 = (-1, a)\)；\(q_1' = (1, -a)\)、\(q_2' = (1, 0)\)、\(q_3' = (1, a)\)。令 K 为集合 \(p_1 + rL\)、\(p_2 + rL\)、\(p_3 + rL\) 与线段 \([q_1, p_1 - rv]\)、\([q_2, p_2 - ru]\)、\([q_3, p_3 - ru]\)、\([q_1', p_1 + ru]\)、\([q_2', p_3 + ru]\)、\([q_3', p_3 + rv]\) 的并。
  \end{enumerate}
\end{enumerate}

令 \((x_{n})_{n\in \mathbf{Z}}\) 是严格递增的实数族，满足 \(\lim_{n\to -\infty}x_n = -1\) 且 \(\lim_{n\to +\infty}x_n = 1\)；对 \(n\in \mathbf{Z}\)，置 \(z_{n} = y_{n} = 1 - |x_{n}|\)。对 \(n\in \mathbf{Z}\)，令 \(f_{n}\) 为从 \([-1,1]^{3}\) 到 \(\mathbf{R}^3\) 的映射，由下式给出：

\[
f _ {n} (x, y, z) = \frac {1 - x}{2} (x _ {n}, y y _ {n}, z z _ {n}) + \frac {1 + x}{2} (x _ {n + 1}, y y _ {n + 1}, z z _ {n + 1})
\]

，并置 \(A_{n}=f_{n}(K)\) 和 \(B_{n}=f_{n}([-1,1]^{3})\)。令 A 为族 \((\mathrm{A}_{n})_{n\in\mathbf{Z}}\) 的并。

\begin{enumerate}
\def\labelenumi{\alph{enumi})}
\item
  证明 A 同胚于 [0, 1]。

(6) 此例出自 E. ARTIN 与 R. FOX 的文章《三维空间中的一些野胞与球面》，载于《数学年刊》49 (1948)，第 979–990 页。

\item
对于每个自然整数 \(m\)，置
\end{enumerate}

\[
\mathrm{A} _ {m} ^ {\prime} = \bigcup_ {n <   - m} \mathrm{B} _ {n} \cup \bigcup_ {- m \leqslant n \leqslant m} \mathrm{A} _ {m} \cup \bigcup_ {n > m} \mathrm{B} _ {n}.
\]

证明 \(R^{3}-A_{m}^{\prime}\) 道路连通，且其基本群由元素 \(a_{n}, b_{n}, c_{n}\)（其中 \(n \in Z\) 且 \(-m \leqslant n \leqslant m\)）生成，并满足关系 \(b_{-m} = c_{-m}a_{-m}\)、\(c_{n+1}a_{n+1} = c_{n}c_{n+1}\)、\(c_{n+1}b_{n} = a_{n}c_{n+1}\)、\(b_{n}c_{n+1} = b_{n+1}b_{n}\)（其中 \(n \in Z\) 且 \(-m \leqslant n < m\)）以及 \(c_{m}a_{m} = b_{m}\)。

\begin{enumerate}
\def\labelenumi{\alph{enumi})}
\setcounter{enumi}{2}
\tightlist
\item
   证明 \(R^{3}-A\) 道路连通，且其基本群由元素 \(a_{n}, b_{n}, c_{n}\)（\(n \in Z\)）生成，并满足关系
\end{enumerate}

\[
b _ {n} = c _ {n} a _ {n}, \quad c _ {n + 1} a _ {n + 1} = c _ {n} c _ {n + 1}, \quad c _ {n + 1} b _ {n} = a _ {n} c _ {n + 1}, \quad b _ {n} c _ {n + 1} = b _ {n + 1} b _ {n}
\]

 当 \(n \in Z\) 时。

d) 证明存在从 \(\pi_1(\mathbf{R}^3 -\mathbf{A})\) 到群 \(\mathfrak{S}_5\) 的同态，它把 \(c_{n}\) 的道路类映为置换 \((1,2,3,4,5)\mapsto (2,3,4,5,1)\)（当 \(n\) 为奇数），并映为置换 \((1,2,3,4,5)\mapsto (4,3,5,2,1)\)（当 \(n\) 为偶数）。推出 \(\mathbf{R}^3 -\mathbf{A}\) 不是道路单连通的。

\begin{enumerate}
\def\labelenumi{\arabic{enumi})}
\setcounter{enumi}{33}
\tightlist
\item
  设 X 是道路连通的拓扑空间，G 是连续左作用于 X 的离散群。令 M 是 X 的开子集，满足 G·M = X。按照 I，第 136 页命题 10 定义集合 S、T、群 F 及群同态 \(\varphi\colon\mathrm{F}\to\mathrm{G}\)。令 \(a\) 是 M 中的一点。
\end{enumerate}

\begin{enumerate}
\def\labelenumi{\alph{enumi})}
\item
  设 \(c\colon\mathbf{I}\to\mathbf{X}\) 是以 \(a\) 为基点的圈。证明存在 F 中唯一的元素 \(g(c)\)，满足以下性质：对每个整数 \(n\) 以及 S 中元素的每个序列 \((s_{1},\ldots,s_{n})\)，若 \(c([(k-1)/n,k/n])\subset s_{k}\mathrm{M}\) 对每个整数 \(k\in\{1,\ldots,n\}\) 都成立，则 \(g(c)=x_{s_{n}^{-1}s_{1}}x_{s_{1}^{-1}s_{2}}\ldots x_{s_{n-1}^{-1}s_{n}}\)。
\item
  证明存在唯一的群同态 \(\gamma\)，从 \(\pi_1(\mathrm{X},a)\) 到 F，使得 \(\gamma ([c]) = g(c)\)，其中 \(c\) 是 X 中以 \(a\) 为基点的每个圈。
\item
  证明 \(\varphi\) 的核等于同态 \(\gamma\) 的像。
\item
  证明 \(\gamma\) 的核由各群 \(\pi_{1}(\mathrm{M}\cup s\cdot\mathrm{M},a)\)（\(s\in S\)）的像的并生成。
\end{enumerate}

\exercisegroup

\begin{enumerate}
\def\labelenumi{\arabic{enumi})}
\tightlist
\item
   设 K 是拓扑域 R、C、H 之一。设 V 是具有可数生成系的 K-向量空间，并赋予其作为拓扑向量空间的最细拓扑（EVT，第 I 章，§1，第 11 页，推论 4）。对每个整数 \(n \geqslant 1\)，令 \(\mathrm{B}_{n}(\mathrm{V})\) 为 \(V^{n}\) 中由自由序列组成的子空间，令 \(\mathrm{G}_{n}(\mathrm{V})\) 为 V 的 n 维子空间所成的集合。记 \(p: \mathrm{B}_{n}(\mathrm{V}) \to \mathrm{G}_{n}(\mathrm{V})\) 为把序列 \(v \in \mathrm{B}_{n}(\mathrm{V})\) 对应到其生成的向量子空间的映射，并在集合 \(\mathrm{G}_{n}(\mathrm{V})\) 上赋予使 p 连续的最粗拓扑。
\end{enumerate}

\begin{enumerate}
\def\labelenumi{\alph{enumi})}
\item
   证明 \(\mathbf{GL}(n,\mathrm{K})\) 在 \(\mathrm{B}_{n}(\mathrm{V})\) 上的右作用可由下式定义：\(v \cdot g = (\sum_{i=1}^{n} v_{i} a_{i,j})_{1 \leqslant j \leqslant n}\)，其中 \(v = (v_{1}, \ldots, v_{n}) \in \mathrm{B}_{n}(\mathrm{V})\)，\(g = (a_{i,j}) \in \mathbf{GL}(n, \mathrm{K})\)。
\item
   证明映射 p 在商空间 \(\mathrm{B}_{n}(\mathrm{V})/\mathbf{G}\mathbf{L}(n,\mathrm{K})\) 到 \(\mathrm{G}_{n}(\mathrm{V})\) 上诱导同胚。推出空间 \(\mathrm{G}_{n}(\mathrm{V})\) 是仿紧的；若 V 有限维，则它是可度量化的。
\item
   证明 \(\mathrm{B}_{n}(\mathrm{V})\) 是以 \(\mathbf{GL}(n,\mathrm{K})\) 为群、以 \(\mathbf{G}_{n}(\mathrm{V})\) 为底的主纤维丛。
\item
   假设 V 是无限维的。证明空间 \(\mathrm{B}_{n}(\mathrm{V})\) 是可缩的。推出空间 \(\mathrm{G}_{n}(\mathrm{V})\) 是群 \(\mathbf{GL}(n,\mathrm{K})\) 的分类空间。
\end{enumerate}

\begin{enumerate}
\def\labelenumi{\arabic{enumi})}
\setcounter{enumi}{1}
\tightlist
\item
  设 H 是无限维的分离预希尔伯特空间。令 \((e_{i})_{i\geqslant0}\) 是生成 H 的稠密子空间的正交归一向量族。对每个整数 \(n\geqslant1\)，令 \(V_{n}(H)\) 为 \(H^{n}\) 中由正交归一序列 \((u_{1},\ldots,u_{n})\) 组成的子空间。
\end{enumerate}

\begin{enumerate}
\def\labelenumi{\alph{enumi})}
\item
  证明存在唯一的连续线性映射 T: H → H，使得对每个整数 i 都有 \(\mathrm{T}(e_{i}) = e_{i+1}\)。证明 \(\|T(x)\| = \|x\|\) 对所有 \(x \in H\) 都成立。
\item
  令 H' 为 H 中与 \(e_{i}\) 正交的向量所成的子空间，其中 i \textless{} n。证明以下条件等价：(i) x 与 \(\mathrm{T}^{n}(x)\) 线性相关；(ii) \(\mathrm{T}(x) = x\)；(iii) \(x \in H'\)。
\item
  证明存在唯一的同伦 \(\sigma\colon\mathrm{V}_{n}(\mathrm{H})\times\mathbf{I}\to\mathrm{V}_{n}(\mathrm{H})\)，使得对每个 \((u_{1},\ldots,u_{n})\in\mathrm{V}_{n}(\mathrm{H})\) 和每个 \(t\in\mathbf{I}\)，\(\sigma((u_{1},\ldots,u_{n}),t)\) 是由族 \((1-t)(u_{1},\ldots,u_{n})+t(\mathrm{T}^{n}(u_{1}),\ldots,\mathrm{T}^{n}(u_{n}))\) 经过正交归一化过程得到的序列。
\item
  证明存在唯一的同伦 \(\tau\colon\mathrm{V}_{n}(\mathrm{H}^{\prime})\times\mathbf{I}\to\mathrm{V}_{n}(\mathrm{H})\)，使得对每个 \((u_{1},\ldots,u_{n})\in\mathrm{V}_{n}(\mathrm{H}^{\prime})\) 和每个 \(t\in\mathbf{I}\)，\(\sigma((u_{1},\ldots,u_{n}),t)\) 是由族 \((1-t)(u_{1},\ldots,u_{n})+t(e_{1},\ldots,e_{n})\) 经过正交归一化过程得到的序列。
\item
  证明空间 \(V_{n}(H)\) 是可缩的。
\end{enumerate}

\begin{enumerate}
\def\labelenumi{\arabic{enumi})}
\setcounter{enumi}{2}
\tightlist
\item
  设 H 是分离的预希尔伯特空间。设 n 是整数且 \(\geqslant 1\)；记 \(V_{n}(H)\) 为由正交归一序列组成的 \(H^{n}\) 的子空间。设 G 是 \(O(n, \mathbf{R})\) 的闭子群。
\end{enumerate}

\begin{enumerate}
\def\labelenumi{\alph{enumi})}
\tightlist
\item
   证明正交群 O(n, R)（LIE，第 IX 章，§3，第 5 号，第 22 页）在空间 \(V_{n}(H)\) 上按下式适当且自由地作用：
\end{enumerate}

\[
(u _ {1}, \dots , u _ {n}) \cdot (a _ {i, j}) = \left(\sum_ {i = 1} ^ {n} a _ {i, j} u _ {i}\right).
\]

\begin{enumerate}
\def\labelenumi{\alph{enumi})}
\setcounter{enumi}{1}
\item
  设 G 是 \(O(n,\mathbf{R})\) 的闭子群。证明典范满射使 \(V_{n}(\mathrm{H})\) 成为以 G 为群、以 \(V_{n}(\mathrm{H})/\mathrm{G}\) 为底的主纤维丛（先处理 \(G = O(n,\mathbf{R})\) 的情形）。
\item
  证明商空间 \(V_{n}(H)/G\) 是可度量化的。（在 \(H^{n}\) 上赋予其自然的预希尔伯特空间结构，并考虑 \(V_{n}(H)\) 上由 \(d(u,v)=\inf_{g\in G}\|u-v\cdot g\|\) 给出的规度 d。）
\item
  证明商空间 \(V_{n}(H)/G\) 是 G 的分类空间。
\item
  推出每个紧李群都有一个作为可度量拓扑空间的分类空间。（应用 LIE，第 IX 章，§9，第 2 号，第 91 页，定理 1。）
\end{enumerate}

\begin{enumerate}
\def\labelenumi{\arabic{enumi})}
\setcounter{enumi}{3}
\item
  通过考察复希尔伯特空间，以类似方式构造酉群 U(n, C)（LIE，第 IX 章，§，第 4 号，第 21 页）的可度量分类空间。
\item
  \begin{enumerate}
  \def\labelenumii{\alph{enumii})}
  \tightlist
  \item
    设 X 是仿紧拓扑空间。证明每个以 X 为底、以 R 为群的主纤维丛都是可平凡化的。
  \end{enumerate}
\end{enumerate}

\begin{enumerate}
\def\labelenumi{\alph{enumi})}
\setcounter{enumi}{1}
\tightlist
\item
  设 X 是亚历山德罗夫半直线（TG，第 IV 章，第 49 页，练习 12）。构造一个以 X 为底、以 R 为群且不可平凡化的主纤维丛。
\end{enumerate}

\begin{enumerate}
\def\labelenumi{\arabic{enumi})}
\setcounter{enumi}{5}
\tightlist
\item
  设 X 为 \(R \times \{0,1\}\) 的商空间，商取使点 \((t,0)\) 与 \((t,1)\) 对每个 \(x \in R_{+}^{*}\) 都等价的最细等价关系。记 p: \(R \times \{0,1\} \to X\) 为典范满射。
\end{enumerate}

\begin{enumerate}
\def\labelenumi{\alph{enumi})}
\item
  证明拓扑空间 X 同伦于一点。
\item
  对于哪些点 \(x \in X\)，带基点空间 (X, x) 是可缩的？
\item
  令 \(U = p(\mathbf{R} \times \{0\})\)，\(V = p(\mathbf{R} \times \{\mathrm{V}\})\)。设 G 是拓扑群。构造从连续映射集合 \(\mathcal{C}(\mathbf{R}_{+}^{*}; \mathbf{G})\)（从 \(R_{+}^{*}\) 到 G）到以 X 为底、以 G 为群的主纤维丛同构类集合的双射。
\item
  推出存在以 X 为底、以 R 为群且不可平凡化的主纤维丛。
\end{enumerate}

\begin{enumerate}
\def\labelenumi{\arabic{enumi})}
\setcounter{enumi}{6}
\tightlist
\item
  设 G 是拓扑群，E 是以 \(S_{2}\) 为底、以 G 为群的主纤维丛。记 \(S_{2}^{+}\) 和 \(S_{2}^{-}\) 为 \(S_{2}\) 与分别由条件 \(z \geqslant 0\) 和 \(z \leqslant 0\) 定义的半空间的交，即两个半球；将 \(S_{2}^{+} \cap S_{2}^{-}\) 识别为 \(S_{1}\)。
\end{enumerate}

\begin{enumerate}
\def\labelenumi{\alph{enumi})}
\item
  证明丛 E 有截面 \(s^{+}\)（分别为 \(s^{-}\)），分别定义在 \(S_{2}^{+}\)（分别在 \(S_{2}^{-}\)）上。证明存在唯一连续映射 \(f\colon S_{1}\to G\)，使得 \(s^{+}(x)=s^{-}(x)\cdot f(x)\) 对所有 \(x\in S_{1}\) 都成立。
\item
  证明从 \(S_{1}\) 到 G 的映射 \(x \mapsto c(x)c(1)^{-1}\) 的严格同伦类不依赖于截面 \(s^{+}\) 和 \(s^{-}\) 的选择；将其记为 \(c(\mathrm{E})\)。
\item
  证明，两个主纤维丛 E 和 \(E'\)（底为 \(S_{2}\)，群为 G）同构，当且仅当 \(c(\mathrm{E}) = c(\mathrm{E}')\)。
\item
  证明对每个类 \(c \in \pi_{1}(G, e)\)，存在以 \(S_{2}\) 为底、以 G 为群的主纤维丛 E，使得 \(c(E) = c\)。
\end{enumerate}

\begin{enumerate}
\def\labelenumi{\arabic{enumi})}
\setcounter{enumi}{7}
\tightlist
\item
  设 X 是仿紧拓扑空间，令 S(X) 为其悬垂（I，第 149 页，练习 1）；记 \(p\colon X \times [0,1] \to S(X)\) 为典范满射。令 \(a\) 为 X 中的一点。
\end{enumerate}

\begin{enumerate}
\def\labelenumi{\alph{enumi})}
\item
  设 G 是拓扑群，E 是以 S(X) 为底、以 G 为群的主纤维丛。证明 E 限制到子空间 \(p(\mathrm{X} \times [0,1/2])\) 后有连续截面 \(s_{0}\)，且限制到子空间 \(p(\mathrm{X} \times [1/2,1])\) 后有连续截面 \(s_{1}\)。
\item
  证明存在唯一映射 \(f\colon \mathrm{X}\to \mathrm{G}\)，使得 \(s_1(x) = s_0(x)\cdot f(x)\) 对所有 \(x\in \mathrm{X}\) 都成立。
\item
  证明在 \([(X,a);(G,e)]\) 中，带基点映射 \(x \mapsto f(x)f(a)^{-1}\) 的严格同伦类不依赖于截面 \(s_{0}\) 和 \(s_{1}\) 的选择；将其记为 \(c(\mathrm{E})\)。
\item
  证明，以 S(X) 为底、以 G 为群的两个主纤维丛 E 和 \(E'\) 同构，当且仅当 \(c(\mathrm{E}) = c(\mathrm{E}')\)。
\item
  证明对每个类 \(c \in [(X, a); (G, e)]\)，存在以 S(X) 为底、以 G 为群的主纤维丛 E，使得 \(c(E) = c\)。
\end{enumerate}

\begin{enumerate}
\def\labelenumi{\arabic{enumi})}
\setcounter{enumi}{8}
\tightlist
\item
  我们将复射影直线 \(\mathbf{P}^{1}(\mathbf{C})\) 识别为 \(C^{2}\) 中的直线集。令 E 为 \(\mathbf{C}^{2} \times \mathbf{P}^{1}(\mathbf{C})\) 的子空间，由点对 \(((u, v), x)\) 组成，这些点满足 \(v \in x\) 且 \(|u|^{2} + |v|^{2} = 1\)。在 E 上赋予 \(S_{1}\) 的作用 \(((u, v), x) \cdot z = ((zu, zv), x)\)。
\end{enumerate}

\begin{enumerate}
\def\labelenumi{\alph{enumi})}
\item
  证明第二投影 \(\mathrm{pr}_{2}\colon\mathrm{E}\to\mathbf{P}^{1}(\mathbf{C})\) 赋予 E 以 \(\mathbf{P}^{1}(\mathbf{C})\) 为底、以 \(S_{1}\) 为群的主纤维丛结构。
\item
  设 \(\varphi\) 是从 \(\mathbf{S}_2\) 到 \(\mathbf{P}^1 (\mathbf{C})\) 的同胚；证明群 \(\pi_1(\mathbf{S}_1,e)\) 由类 \(c(\varphi^{*}\mathrm{E})\) 生成。
\end{enumerate}

\backmatter

\specialchapter{符号索引}

 \(\mathcal{C}_{\mathrm{B}}(\mathrm{X};\mathrm{X}^{\prime})\) (B-态射) \ldots.. 2 \(\mathrm{Isom}_{\mathrm{B}}(\mathrm{X};\mathrm{X}^{\prime})\) (B-同构) \ldots.. 2 \(\mathrm{Aut}_{\mathrm{B}}(\mathrm{X})\) (B-自同构) .. 2 \(\mathrm{X}_{b}\) (纤维) \ldots.. 2 \(\mathrm{X}_{\mathrm{A}}\) (诱导空间) \ldots.. 2 \(\mathrm{X} \times_{\mathrm{B}} \mathrm{X}^{\prime}\) (纤维积) \ldots.. 3 \(\Delta_{\mathrm{X}}\) (对角线) \ldots.. 4 \(\delta_{\mathrm{X}}\) (对角态射) \ldots.. 4 \(f^{*}(\mathrm{X})\) (换基) .. 5 \(\prod_{\mathrm{B}} \mathrm{X}_{i}\) (纤维积) \ldots.. 5 \(\mathcal{C}_{c}(\mathrm{U};\mathrm{V})\) (从 U 到 V 的连续映射) \ldots.. 19 \(\mathcal{S}(\mathrm{A};\mathrm{E})\) (连续截面) .. 33 \(s \mid \mathrm{U}\) (限制) \ldots.. 43 \(\underline{\mathcal{F}}(\mathrm{B};\mathrm{X})\) (映射层) \ldots.. 45 \(\underline{\mathcal{C}}(\mathrm{B};\mathrm{X})\) (连续映射层) \ldots.. 45 \(\underline{\mathcal{L}}(\mathrm{B};\mathrm{E}), \underline{\mathcal{L}}(\mathrm{E})\) (截面层) \ldots.. 45 \(\underline{\mathcal{C}}_{\mathrm{B}}(\mathrm{E};\mathrm{E}^{\prime})\) (B-态射层) \ldots.. 45 \(\underline{\mathcal{Mor}}_{\mathrm{B}}(\mathrm{E};\mathrm{E}^{\prime})\) (B-态射层) \ldots.. 46 \(\underline{\mathcal{Isom}}_{\mathrm{B}}(\mathrm{E};\mathrm{E}^{\prime})\) (同构层) \ldots.. 46 \(\underline{\mathcal{C}}^{r}(\mathrm{X};\mathrm{Y})\) (\(C^{r}\) 类映射的层) \ldots.. 46 \(\underline{\mathfrak{P}}(\mathrm{B})\) (子空间层) \ldots.. 46 \(\mathrm{E}_{\mathcal{F}}\) (平展空间) \ldots.. 49 \(\widetilde{\mathcal{F}}\) (预层的伴随层) 53 \(u_{*}(\mathcal{F})\) (预层、层的直像) 57 \(u^{*}(\mathcal{G})\) (预层、层的逆像) 58 \(\mathcal{G}_{\mathrm{A}}\) (诱导层) 59 \(\mathrm{G}^{\circ}\) (反群) 91 \(\mathcal{C}_{\mathrm{B}}^{\mathrm{G}}(\mathrm{E};\mathrm{E}^{\prime})\) (主纤维丛的态射) 92 \(\mathrm{E} \times^{\mathrm{G}} \mathrm{F}\) (相伴纤维丛) . 106 \(\mathrm{Z}_{\mathrm{cont}}^{1}(\mathrm{B},\mathcal{U},\mathrm{G})\) (连续 1-上闭链) 114 \(\mathrm{H}_{\mathrm{cont}}^{1}(\mathrm{B},\mathcal{U},\mathrm{G})\) (上同调类) 116 \(\mathrm{P}(\mathrm{B},\mathrm{G})\) (主丛的同构类) 118 \(\mathrm{P}(\mathrm{B},\mathcal{U},\mathrm{G})\) (主丛的同构类) 118 \(\mathrm{H}_{\mathrm{cont}}^{1}(\mathrm{B},\mathrm{G})\) (上同调类) 119 \(\mathcal{C}((\mathrm{X},x);(\mathrm{Y},y))\) (带基点连续映射) 120 \(p_{\mathrm{A}}\) (凸部分 A 的规度) 123 \(\mathrm{S}(\mathrm{X})\) (空间 X 的悬垂) 149\\
Som(C) (箭图的顶点) 152\\
Fl(C) (箭图的箭) . 152

 \(o_{C}\) (箭图的源映射) \ldots.. 152 \(t_{C}\) (箭图的靶映射) \ldots.. 152 \(\mathrm{Fl}_{a,b}(\mathrm{C}), \mathrm{C}_{a,b}\) (连接 a 与 b 的箭) \ldots.. 152 Som(φ) (箭图态射) \ldots.. 152 \(\mathrm{Fl}(\varphi)\) (箭图态射) \ldots.. 152 \(c*d\) (拼接道路) \ldots.. 154 \(\pi_{0}(\mathrm{C})\) (连通分支) \ldots.. 155 \(\pi_{0}(\varphi)\) (在连通分支上诱导的映射) .. 155 \(\overline{f}\) (Fl(C) 上的对合) \ldots.. 155 Ch(C) (箭图中的道路集) \ldots.. 160 \(\Omega_{C}\) (箭图的道路范畴) \ldots.. 160 \(f^{-1}\) (可逆箭的逆) \ldots.. 161 \(G_{a}\) (迷向群) \ldots.. 162 Int(f) (迷向群的同构) \ldots.. 162 \(\varphi_{a}\) (迷向群的同态) \ldots.. 162 Orb(G) (群胚的轨道) \ldots.. 162 Orb(φ) (通过取轨道由 φ 诱导的映射) . 162 \(\mathcal{B}(X,p)\) (置换群胚) \ldots.. 163 \(\varinjlim_{i}G_{i}\) (群胚的归纳极限) \ldots.. 164 X ×s X (映射的配对群胚) \ldots.. 165 \(\varphi^{*}G\) (逆像群胚) \ldots.. 166 Ker(φ) (群胚态射的核) \ldots.. 166 G/H (商群胚) \ldots{} 170 \(\overline{\varphi}\) (通过取商诱导的群胚态射) . 171 \(\varpi_{G}\) (道路类群胚) \ldots.. 172 \(\varpi(\varphi)\) (道路类群胚态射) \ldots.. 173 F(A) (自由群) \ldots.. 173 Grp(C) (自由群胚) \ldots. 174

G/F (由箭缩并得到的群胚) \ldots{} 176 π₁(G, a) (图在 a 处的庞加莱群) \ldots{} 178 u ∼ v (群胚态射的同伦) \ldots{} 181 Coh(φ, ψ) (上同伦子) \ldots{} 184 * Gᵢ (群的自由积) i∈I 193 Coeg(φ, ψ) (余均衡子) \ldots{} 199 C(X; Y) (从 X 到 Y 的连续映射) \ldots{} 230 {[}X; Y{]} (连续映射的同伦类) . 230 f (连续映射的同伦类) \ldots{} 230 C((X, x); (Y, y)) (带基点连续映射) \ldots{} 232 Sₙ (欧氏球面) \ldots{} 234 Cyl(f) (映射的柱) \ldots{} 237 X ∪ f B (附着所得拓扑空间) \ldots{} 248 sₓ/A (X/A 的基点) \ldots{} 252 Côn(f) (映射的锥) 253 c * d (拼接道路) \ldots{} 256 ¯ c (逆道路) \ldots{} 256 Λ(X) (道路空间) \ldots{} 257 Λₓ(X) (起点为 x 的道路) . 257 λₓ,y(X) (起点为 x、终点为 y 的道路) \ldots{} 257 π₀(X) (道路连通分支) \ldots{} 259 π₀(f) (在道路连通分支上由 f 诱导的映射) \ldots{} 259 ϖₓ,y(X) (从 x 到 y 的道路类) \ldots{} 290 Ω(X) (环路空间) \ldots{} 290 Ωₓ(X) (x 处的圈) \ldots{} 290 εₓ (常值圈) \ldots{} 290 ϖ(X) (庞加莱群胚) 292 π₁(X, x) (基本群) 292 ϖ(f) (f 诱导的庞加莱群胚态射) \ldots{} 294 f*(γ) (ϖₓ,y(f)(γ) 的缩写) \ldots{} 294

 \(x \cdot \gamma\) (庞加莱群胚的作用) \ldots.. 304 \(e_{\mathrm{B}}\) (靶映射) \ldots.. 342 \(\varepsilon_{\mathrm{B}}\) (靶映射) \ldots.. 342 \(\pi_{1}(\mathrm{G})\) (拓扑群的庞加莱群) \ldots.. 375 \(\mathrm{R}_{\tau} \{\xi z_{1}, z_{2}\}\) (与下降数据相关联的等价关系) \ldots.. 383

\((Z/R_{\tau},q)\) (由下降数据商得的空间 \((Z,p)\)) \ldots. 384 \(C_{\tau,\tau'}(Z;Z')\) (与下降数据相容的态射) \ldots. 387 \textbar G\textbar{} (图的几何实现) \ldots. 417 \(\bigvee_{i}(X_{i},x_{i})\) (带基点拓扑空间的楔和) \ldots{} 421

\specialchapter{术语索引}

\indexletter{A}

伴随, 59 容许作用 ---, 317 子群 ---, 315 拓扑 ---, 316 道路连续映射, 267 平展, 28 柔软, 64 带基点, 120 相对连通, 353 分离, 25 严格, 18 拓扑平展, 28 普遍严格, 20 树, 157, 174 极大, 158 极大定向, 178 边, 定向边, 155 框架, 185, 201, 395, 406 覆叠的, 409 附着 附着（拓扑空间由 --- 得到）, 248

\indexletter{B}

锥的底, 253, 426 柱的底, 238 夏威夷耳环, 146

带基点拓扑空间的楔和, 421 靶 道路的终点, 154 箭的靶, 151

\indexletter{C}

箭图, 151 与定向图相关的, 156 --- 的连通分支, 155 连通的, 155 凯莱, 221 施赖尔, 221 反, 216 --- 的积, 153 方形, 7 笛卡尔, 7 复合, 15 纤维, 4 范畴, 160 箭图的道路范畴, 160 反, 224 凯莱 凯莱（--- 的图）, 221 B-空间的换基, 4 主丛的, 93 道路, 160 箭图中的, 154 拓扑空间中的, 256 拼接的, 154, 256

箭图中---的长度, 154 逆, 157, 256 连接 x 与 y, 256 无折返, 157, 173 道路 道路空间的---, 257 --- 类的逆, 290 可拼接的, 154, 256 --- 的提升性质, 302 严格同伦的, 289 图中的道路类, 172 上同调类, 119 分类（--- 空间）, 448 上闭链, 114 --- 的上同调类, 119 连续的, 114 由一族平凡化定义的, 115 平凡的, 118 上闭链 同调的, 115 余维, 358 余均衡子, 199 上同伦子, 185, 199 道路连通分支, 259 可拼接道路类的复合, 290 锥 ---的底, 253, 426 拓扑空间的锥, 254 映射的锥, 253, 425 拓扑空间的开锥, 254 ---的顶点, 254, 425 连通 相对---映射, 353 道路连通拓扑空间, 258 局部道路连通拓扑空间, 260 单连通拓扑空间, 124 道路单连通拓扑空间, 340 缩并的 拓扑空间, 236 可缩的

拓扑空间 一点处---, 234 局部---拓扑空间, 346 带基点---拓扑空间, 234 缩并, 236 锥到其底的典范缩并, 254 柱到其底的典范缩并, 239 缩并所得拓扑空间, 252 强缩并, 237 凸, 122 度量空间的筛分, 362 柱 ---的底, 238 映射的柱, 237

\indexletter{D}

覆叠的次数, 73 可解圈（拓扑空间 ---）, 340 庞加莱上半平面, 150 对角线, 4 下降数据, 382 典范, 383 有效, 384 与 --- 相容的态射, 387 范坎彭数据, 406, 433

\indexletter{E}

有效（下降数据 ---）, 384 均衡子, 153, 165 单位元, 160 康托尔集, 270 有限特征子集的, 279 带基点, 120 右滤预序, 164 无间隙全序, 274 邻偶, 272 补结构, 208, 212, 407 完备, 208, 407 基本, 192, 194, 206, 211, 407 空间

分类, 448 霍普夫, 373 道路的, 257 圈的, 290 平展, 29 纤维丛 可伴随于主纤维丛的, 107 与主纤维丛相伴的, 106 局部平凡, 68 主, 91 由上闭链定义的主, 115 平凡主, 92 万有主, 448 平凡, 68 可平凡化, 68 拓扑空间 道路连通, 258 贝尔, 362 可解圈, 340 局部道路连通, 260 局部可缩, 346 可度量化, 267 附着所得, 248 缩并所得, 252 识别所得, 430 带基点, 232 波兰, 364 积, 297 单连通, 124 道路单连通, 340 苏斯林, 362 B-空间, 1 与预层相伴的平展, 50 柔软平展, 64 诱导, 2 积, 5 商, 3 星形 星形（ 一点处的---集）, 234

\indexletter{F}

层, 43 与预层相伴的, 53 常值, 86 映射的, 45 连续映射的, 45

C 类映射 \(^{r}\) , 46 B-态射的, 45 连续截面的, 45 子空间的, 46 由限制到开集得到的, 44 局部常值映射的, 45 主纤维丛态射的, 446 B-空间同构的, 46 终端, 142 松弛, 144 局部常值, 86 柔软, 64 积, 46

层，---的同构, 47；---的态射, 47。

局部平凡纤维化, 68 主的, 91

纤维, 2

固定子, 166

箭 ---的靶, 151 箭图的, 151 可逆的, 161 逆的, 155 ---的起点, 151 ---的源, 151 ---的端点, 151

可复合的箭, 159 一族可复合的---, 160 一族可复合---的积, 160

函子, 161

森林, 157, 178 极大, 157 极大定向, 183, 188

\indexletter{G}

芽, 50 图, 155 --- 完备, 220 与箭图相关的, 156 组合的, 158 凯莱, 221

施赖尔的, 221 图的定向, 156 ---的几何实现, 417 有向图, 156 与箭图相关的有向图, 156 有限表示群, 359 迷向群, 162 基本群, 292 自由群, 173, 193 庞加莱群, 292 图在顶点处的, 178 拓扑群的, 375 群胚, 162 与群相伴的, 163 余均衡子, 394 庞加莱的, 292 图的道路类的, 172 配对的, 163, 165 置换的, 163 基本群胚, 205, 292 逆像, 166, 183 自由群胚, 174 由箭缩并得到的, 176 庞加莱群胚的典范作用, 304 群胚的作用, 167 积, 164, 297 商, 170 单传递, 162, 174 传递, 162, 168, 170, 192, 197, 395 群胚 一族---的归纳极限, 164

\indexletter{H}

同伦的空间对 ---, 235 同伦的拓扑空间 ---, 233 同伦, 232 空间对的, 234 局部同态, 369 连续同态, 369 同伦的

群胚的态射 ---, 180 同伦, 229 与上同伦子相关的典范态射, 185 ---的复合类, 231 ---的可逆类, 232 ---的带基点类, 232 连续映射的---类, 230 模---的单位元, 373 在集上固定, 229 模---的逆, 233 ---类的逆, 233 拼接的, 230 自由的, 299 ---的起点, 229 带基点的, 231 模---的互逆, 233 ---关系, 230 带基点的---关系, 232 严格的, 289, 340 ---的端点, 229 可拼接同伦, 229 ---的延拓性质, 240 同伦等价态射, 181 霍普夫（---空间）, 373

\indexletter{I}

模同伦的逆, 181 道路类的逆, 290 可逆箭的逆, 161 B-空间的同构, 1 范畴的同构, 161 层的同构, 47 预层的同构, 47 局部同构, 369 B-同构, 1

\indexletter{J}

规度, 123 可拼接道路 ---, 154, 256 可拼接同伦 ---, 229 道路序列 ---, 291 道路的拼接, 154, 160

\indexletter{L}

常值圈, 154, 340 箭图中的, 154 拓扑空间中的, 256, 303 一点 x 处的, 290 圈 一点 x 处---的类, 290 ---的空间, 290 自由同伦的, 299 列车的宽度, 264 箭图中的合成法则, 159 结合的, 160 列车的长度, 264

\indexletter{M}

与余均衡子相关的典范态射,

与上同伦子相关的典范态射, 185 与下降数据相容的, 387 B-空间的, 1 箭图的, 152 范畴的, 161 层的, 47 连续映射诱导的庞加莱群胚的,

通过取商由态射得到的群胚的,

延拓箭图态射的群胚的, 175 预层的, 47 投影的, 48 主纤维丛的, 91 对角的, 4 拓扑群的局部, 369 B-态射, 1 由基本配备诱导的态射,

主纤维丛的 f-态射, 94

N 尼尔森–施赖尔（---定理）, 418 核（群胚态射的）, 166 环面结, 468

\indexletter{O}

对象, 161 基本群的容许作用,

基本群的典范作用,

庞加莱群胚的典范作用, 304 群胚的, 167, 313, 383 群的自由作用, 213 逆紧, 287, 349 庞加莱群胚的无单值, 313 群胚的无单值, 168 庞加莱群胚的无局部单值, 313 群胚的轨道, 162 群胚在集上的作用的, 168 定向, 156 图的边的, 155 起点映射 ---, 257 道路的, 154, 256 箭的, 151 同伦的, 229

\indexletter{P}

可达部分, 361 相对于一对的特选部分, 197 贫部分, 360 PCV（性质 ---）, 37 全子箭图, 152 子群胚, 293 基点, 120, 252 预层, 42 由限制到开集得到的, 43

直像, 57 逆像, 58 分离, 143 预层 ---的同构, 47 ---的态射, 47 积 层的, 46 主纤维丛的, 93 丛, 3 主纤维丛的丛, 93 群的自由积, 193 投影, 1 指标 j 的, 5 性质 (PCV), 37 (VK), 405 弗拉格门–布劳威尔的, 467 道路提升的, 302, 303, 394 同伦扩张的, 240 群胚归纳极限的泛性质, 164 柱的, 238 自由积的, 193, 209 余均衡子的, 199 上同伦子的, 185, 187 群胚积的, 164 锥的, 254 自由群的, 175 箭图积的, 153 连续统的基数, 366

\indexletter{R}

列车的加细, 264 图的几何实现,

关系元, 372 与下降数据相关联的等价关系, 383 连续提升, 2, 33 道路的, 302 道路---的性质, 302 层截面的限制, 43 收缩核

绝对, 322 绝对邻域, 322 拓扑空间的, 235 收缩映射, 235 去掉锥顶点的锥到其底的典范收缩, 254 柱到其底的典范收缩, 238 覆叠, 69 可伴随于主覆叠的, 108 有限, 75 伽罗瓦, 101, 312 局部有限, 75 主, 97 积, 88 万有, 120, 343

\indexletter{S}

施赖尔 施赖尔（---的图）, 221 连续截面, 2, 33 （预）层的, 43 线段, 122 开的, 122 分离（映射 ---）, 25 单传递（群胚 ---）, 162 一族 B-空间的和, 2 顶点 箭图的, 151 锥的, 254, 255, 425 源 道路的, 154 箭的, 151 子箭图, 152 全的, 152 子范畴, 164 子层, 44 子图, 155 基本群的容许子群, 315 子群胚, 165 正规的, 168 由箭集生成的正规, 169 由子箭图生成的, 165

由定向子图生成的, 292 全的, 193, 293 子预层, 44 球面, 234 严格映射 ---, 18 普遍严格映射 ---, 20 一致结构, 272 层截面的支集,

拓扑空间的悬垂,

\indexletter{T}

终点映射 ---, 257 道路的终点, 154, 256 箭的终点, 151 同伦的终点, 229 哈恩–马祖尔凯维奇定理, 272 若尔当定理, 468 尼尔森–施赖尔定理, 418 谢拉赫定理, 366 范坎彭定理, 431

茎, 50 容许拓扑, 316 列车, 264 ---的宽度, 264 ---的长度, 264 ---的加细, 264 ---的车厢, 264 传递（---群胚）, 162 结构的传输, 166 可平凡化的---纤维空间, 68 ---主纤维空间, 92 典范平凡化, 95 纤维空间的, 68 主纤维空间的, 92

\indexletter{V}

范坎彭数据 ---的, 406, 433 由---的数据得到的完备配备, 407 ---的定理, 431

W W 列车的车厢, 264
